% ============================================================
\section{Characteristic Realization}
\label{sec:characteristic-realization}
% ============================================================

The calculus of Part~I admits a natural realization on conservative systems
that can be organized simultaneously by observer cuts, transported
characteristics, and variational trajectories. In the language introduced
earlier, this is the passage from the abstract horizon complex to a realized
characteristic field on that complex. Along characteristics, conservative
transport collapses bulk evolution to boundary flux. On a horizon complex,
finite horizon cuts reproduce the exact defect algebra of Part~I. When the
same dynamics arise from a regular action principle, the observer-side
characteristic field is recovered from the classical variational triad.
Together these constructions furnish a bridge from conservative transport to
the operator calculus. The transport-side Eulerian/Lagrangian/characteristic
descriptions and the dynamics-side Lagrangian/Hamiltonian/Hamilton--Jacobi
descriptions are classical; see, for example,
\cite{Evans2010,Arnold1989}.

In this realized setting, the abstract vocabulary of the closure-balance
realization acquires the standard dynamical labels used in applications. A
realized primitive will be called a \emph{characteristic}, and a realized
incidence of such primitives will be called an \emph{event}. Nothing in the
structural calculus depends on those labels; they serve only to place the
realized chapter in standard dynamical language.

\subsection{Conservative Characteristic Transport}

Let $J \subset \bbR$ be a time interval, let $M$ be a smooth manifold with
volume form $\mu$, and let $V_t$ be a smooth time-dependent vector field on
$M$. Write $\Phi_{t,s}$ for the associated flow, and define the spacetime
manifold $X := J \times M$ with spacetime vector field
\[
\widetilde V := \partial_t + V_t
\]
and spacetime volume form
\[
\nu := dt \wedge \mu.
\]
Assume throughout that the flow is conservative in the sense
\[
\mathcal L_{V_t}\mu = 0 \qquad \text{for all } t \in J.
\]
Then $\mathcal L_{\widetilde V}\nu = 0$.

\begin{definition}[Eulerian Presentation]
\label{def:eulerian-presentation}
\lean{EulerianPresentation}
For a time-dependent scalar observable $\Theta_t \in C^\infty(M)$, the
\emph{Eulerian presentation} of conservative transport is the first-order
equation
\[
\mathcal D_\chi \Theta := \partial_t \Theta + \mathcal L_{V_t}\Theta = 0.
\]
\end{definition}

\begin{definition}[Lagrangian Presentation]
\label{def:lagrangian-presentation}
\lean{LagrangianPresentation}
The \emph{Lagrangian presentation} of the same transport law is the constancy of
the pullback along the flow:
\[
\frac{d}{dt}\bigl(\Phi_{t,0}^* \Theta_t\bigr)=0.
\]
\end{definition}

\begin{definition}[Characteristic Presentation]
\label{def:characteristic-presentation}
\lean{CharacteristicPresentation}
Fix $a \in M$ and let $\gamma_a(t):=\Phi_{t,0}(a)$ be the characteristic
through $a$. The \emph{characteristic presentation} is the constancy of the
transported quantity along each such curve:
\[
\frac{d}{dt}\Theta_t(\gamma_a(t))=0.
\]
\end{definition}

\begin{figure}[ht]
\centering
\begin{tikzpicture}[node distance=3.2cm, scale=1.05, every node/.style={scale=1.05}]
  \node[ccbox, align=center, minimum width=2.8cm] (E) at (-3.2,0) {\textbf{Eulerian}\\$\partial_t \Theta + \mathcal L_{V_t}\Theta = 0$};
  \node[ccbox, align=center, minimum width=2.8cm] (L) at (3.2,0) {\textbf{Lagrangian}\\$\dfrac{d}{dt}(\Phi_{t,0}^*\Theta_t)=0$};
  \node[ccbox, align=center, minimum width=2.8cm] (C) at (0,2.8) {\textbf{Characteristic}\\$\dfrac{d}{dt}\Theta_t(\gamma_a(t))=0$};

  \draw[ccarrow] (E) -- node[below, cclabel] {pullback by $\Phi_{t,0}$} (L);
  \draw[ccarrow] (E) -- node[left, cclabel] {restriction to $\gamma_a$} (C);
  \draw[ccarrow] (L) -- node[right, cclabel] {evaluation on labels $a$} (C);
\end{tikzpicture}
\caption{Observer-side compatibility of the Eulerian, Lagrangian, and characteristic presentations of a conservative transport law.}
\label{fig:observer-triad}
\end{figure}

\begin{proposition}[Observer-Side Compatibility]
\label{prop:observer-triad}
\lean{observer_triad}
Under the standing hypotheses above, Definitions~\ref{def:eulerian-presentation},
\ref{def:lagrangian-presentation}, and
\ref{def:characteristic-presentation} are equivalent.
\end{proposition}

(Proof in Appendix~\ref{sec:appendix-dynamics}.)

\begin{definition}[Characteristic Flux Form]
\label{def:characteristic-flux-form}
\lean{CharacteristicFluxForm}
For a transported scalar $\Theta$ on $X$, define the \emph{characteristic flux
form}
\[
J_\Theta := \Theta \,\iota_{\widetilde V}\nu.
\]
\end{definition}

\begin{lemma}[Flux Differential Identity]
\label{lem:flux-differential}
\lean{flux_differential}
With notation as above,
\[
dJ_\Theta = (\mathcal D_\chi \Theta)\,\nu.
\]
\end{lemma}

(Proof in Appendix~\ref{sec:appendix-dynamics}.)

\begin{theorem}[Cut Conservation]
\label{thm:cut-conservation}
\lean{cut_conservation}
Let $U \subset X$ be a compact spacetime tube whose boundary decomposes as
\[
\partial U = \Sigma_1 - \Sigma_0 + \Sigma_{\mathrm{side}},
\]
where $\Sigma_0$ and $\Sigma_1$ are transverse cuts and
$\widetilde V$ is tangent to the side boundary
$\Sigma_{\mathrm{side}}$. If $\mathcal D_\chi \Theta = 0$ on $U$, then
\[
\int_{\Sigma_1} J_\Theta = \int_{\Sigma_0} J_\Theta.
\]
\end{theorem}

(Proof in Appendix~\ref{sec:appendix-dynamics}.)

\begin{theorem}[Characteristic Integration by Parts]
\label{thm:characteristic-ibp}
\lean{characteristic_ibp}
Let $f,g \in C^\infty(U)$ on a compact spacetime tube $U \subset X$. Then
\[
\int_U \bigl((\mathcal D_\chi f)g + f(\mathcal D_\chi g)\bigr)\,\nu
=
\int_{\partial U} fg\,\iota_{\widetilde V}\nu.
\]
\end{theorem}

(Proof in Appendix~\ref{sec:appendix-dynamics}.)

\begin{corollary}[Characteristic Boundary Reduction]
\label{cor:characteristic-boundary-reduction}
\lean{characteristic_boundary_reduction}
If $\mathcal D_\chi f = 0$ on $U$, then
\[
\int_U f(\mathcal D_\chi g)\,\nu
=
\int_{\partial U} fg\,\iota_{\widetilde V}\nu.
\]
Thus conservative transport along characteristic directions removes the bulk
term and leaves only the boundary contribution.
\end{corollary}

\begin{proof}
Substitute $\mathcal D_\chi f = 0$ into
Theorem~\ref{thm:characteristic-ibp}.
\end{proof}

The same elimination mechanism has a direct cellular formulation closer to the
finite observer setting of the structural calculus.

\begin{definition}[Cellular Conservative Current]
\label{def:cellular-current}
\lean{CellularConservativeCurrent}
Let $K$ be a locally finite oriented cell complex. A cochain
$J \in C^1(K;\bbR)$ is called a \emph{cellular conservative current} on a finite
$2$-chain $C$ if
\[
\langle \delta J,\sigma\rangle = 0
\]
for every $2$-cell $\sigma$ in the support of $C$.
\end{definition}

\begin{proposition}[Cellular Cut Conservation]
\label{prop:cellular-cut-conservation}
\lean{cellular_cut_conservation}
Let $\Gamma_0,\Gamma_1$ be finite $1$-cycles in $K$ such that
\[
\Gamma_1-\Gamma_0=\partial C
\]
for some finite $2$-chain $C$. If $J$ is a cellular conservative current on
$C$, then
\[
\langle J,\Gamma_1\rangle = \langle J,\Gamma_0\rangle.
\]
\end{proposition}

(Proof in Appendix~\ref{sec:appendix-dynamics}.)

Proposition~\ref{prop:cellular-cut-conservation} is the discrete companion of
Theorem~\ref{thm:cut-conservation}: in both regularity classes, conservative
transport reduces finite observation to boundary data on cuts.

\subsection{Horizon Realization of the Calculus}

Characteristic transport supplies the boundary law. To recover the operator
calculus of Part~I, the full state space must also carry a closed differential
together with a finite horizon decomposition.

\begin{definition}[Characteristic Horizon Realization]
\label{def:characteristic-horizon-realization}
\lean{CharacteristicHorizonRealization}
A \emph{characteristic horizon realization} consists of the following data.
\begin{enumerate}[label=(\roman*), leftmargin=2em, itemsep=0.2em]
\item A Hilbert complex $(\cH^\bullet,d)$ with $d^2=0$.
\item A nested family of orthogonal projections
      $\{\CCPi{h}\}_{h \in \bbN}$ on each degree of $\cH^\bullet$.
\item A strongly continuous unitary transport family $\{U_t\}_{t \in J}$ on
      $\cH^\bullet$ such that
      \[
      U_t d = d U_t,
      \qquad
      U_t \CCPi{h} = \CCPi{h} U_t
      \quad \text{for all } h \in \bbN.
      \]
\end{enumerate}
The realized observed operator, leakage operator, and defect operator are
defined exactly as in Part~I:
\[
\CCObs{d}{h} = \CCPi{h} d \CCPi{h},
\qquad
\CCLeak{d}{h} = \CCPibar{h} d \CCPi{h},
\qquad
\CCDefect{d}{h} = \CCPi{h} d \CCPibar{h} d \CCPi{h}.
\]
\end{definition}

\begin{figure}[ht]
\centering
\begin{tikzpicture}[node distance=2.8cm, scale=1.05, every node/.style={scale=1.05}]
  \node[ccbox, align=center, minimum width=2.6cm] (bulk) at (0,2.8) {\textbf{Full Horizon Complex}\\$d^2 = 0$};
  \node[ccbox, align=center, minimum width=2.6cm] (obs) at (-3,0) {\textbf{Observable Sector}\\$\CCObs{d}{h}=\CCPi{h} d \CCPi{h}$};
  \node[ccbox, align=center, minimum width=2.6cm] (sh) at (3,0) {\textbf{Shadow Sector}\\$\CCPibar{h}\cH$};
  \node[ccbox, align=center, minimum width=3.0cm] (def) at (0,-2.6) {\textbf{Observed Defect}\\$\CCObs{d}{h}^2 = -\CCDefect{d}{h}$};

  \draw[ccarrow] (bulk) -- node[left, cclabel] {$\CCPi{h}$} (obs);
  \draw[ccarrow] (bulk) -- node[right, cclabel] {$\CCPibar{h}$} (sh);
  \draw[ccarrow] (obs) -- node[above, cclabel] {$\CCLeak{d}{h}$} (sh);
  \draw[ccarrow] (sh) -- node[below, cclabel] {$\CCPi{h} d$} (obs);
  \draw[ccarrow] (obs) -- node[right, cclabel] {round trip} (def);
\end{tikzpicture}
\caption{Full closure holds on the horizon complex. Finite observation inserts a horizon cut, and the observed failure of closure is exactly the round trip through shadow.}
\label{fig:bulk-to-defect}
\end{figure}

\begin{lemma}[Through-Horizon Decomposition]
\label{lem:through-horizon-decomposition}
\lean{through_horizon_decomposition}
For every observed state $u \in \CCPi{h}\cH^\bullet$,
\[
du = \CCObs{d}{h}u + \CCLeak{d}{h}u.
\]
Applying $d$ once more and projecting back to the observed sector gives
\[
\CCPi{h}d\,\CCObs{d}{h}u = -\,\CCPi{h}d\,\CCLeak{d}{h}u = -\,\CCDefect{d}{h}u.
\]
\end{lemma}

(Proof in Appendix~\ref{sec:appendix-dynamics}.)

\begin{theorem}[Characteristic Realization of HFT]
\label{thm:characteristic-realization}
\lean{characteristic_realization}
Let $(\cH^\bullet,d,\{\CCPi{h}\},\{U_t\})$ be a characteristic horizon
realization. Then, for every horizon $h$,
\[
\CCObs{d}{h}^2 = - \CCDefect{d}{h}.
\]
Moreover,
\[
U_t \CCObs{d}{h} = \CCObs{d}{h} U_t,
\qquad
U_t \CCDefect{d}{h} = \CCDefect{d}{h} U_t,
\]
so the realized observed differential and its defect are transported
equivariantly by the conservative evolution.
\end{theorem}

(Proof in Appendix~\ref{sec:appendix-dynamics}.)

\begin{theorem}[Characteristic horizon realizations transport the horizon complex]
\label{thm:characteristic-horizon-realizations-transport-the-horizon-complex}\lean{CharacteristicHorizonRealization.horizonComplex}\lean{CharacteristicHorizonRealization.horizonComplexSurface}
Every characteristic horizon realization carries the Part~I horizon-complex
laws unchanged. Concretely, on the same realized carrier:
\begin{enumerate}[label=(\roman*), leftmargin=2em, itemsep=0.2em]
\item the full boundary still satisfies \(d^2=0\);
\item on every cut, the observed boundary still satisfies
      \(\CCObs{d}{h}^2=-\,\CCDefect{d}{h}\);
\item the conservative transport intertwines both the observed boundary and the
      defect:
      \[
      U_t\CCObs{d}{h}=\CCObs{d}{h}U_t,
      \qquad
      U_t\CCDefect{d}{h}=\CCDefect{d}{h}U_t.
      \]
\end{enumerate}
So a characteristic horizon realization is best understood as a transport
realization of a horizon complex, not merely as an external flow laid on top of
the Part~I state space.
\end{theorem}

\begin{proof}
The associated horizon complex uses the same tower and the same full boundary
\(d\). The closure and cut-defect clauses are exactly the transported Part~I
laws, and the transport identities are the content of
Theorem~\ref{thm:characteristic-realization}.
\end{proof}

\begin{corollary}[Transfer to the Part~I Calculus]
\label{cor:transfer-part-one}
\lean{transfer_part_one}
Every statement of Part~I whose hypotheses involve only the closed operator
$d$, the horizon tower $\{\CCPi{h}\}$, and the algebra they generate applies to
every characteristic horizon realization.
\end{corollary}

\begin{proof}
By Definition~\ref{def:characteristic-horizon-realization}, the realized data
has exactly the same formal type as the operator interface of Part~I. By
Theorem~\ref{thm:characteristic-realization}, the realized observed operator
satisfies the Horizon Fundamental Theorem with the same proof as in the abstract
setting. Any further theorem in Part~I that depends only on that interface
therefore transports unchanged to the realized class.
\end{proof}

\subsection{Variational Realization and the Dynamics-Side Triad}

The observer-side triad concerns the transport of conserved quantities along a
given characteristic field. A different question concerns the origin of that
field when the dynamics are generated by an action principle. The corresponding
classical triad is Lagrangian, Hamiltonian, and characteristic. Here the
Hamiltonian formulation is the phase-space evolution law, whereas the
characteristic formulation refers to the characteristic curves of the
Hamilton--Jacobi equation. The distinguished trajectories are the stationary
action curves seen simultaneously in both languages
\cite{Arnold1989}.

\begin{definition}[Regular Variational System]
\label{def:regular-variational-system}
\lean{RegularVariationalSystem}
Let $Q$ be a smooth configuration manifold and let
$L \in C^2(TQ \times J)$ be a Lagrangian whose fiber Hessian in the velocity
variables is everywhere invertible. The Legendre transform
\[
\mathbb F L : TQ \to T^*Q,
\qquad
(q,\dot q) \mapsto \Bigl(q,\frac{\partial L}{\partial \dot q}\Bigr)
\]
then defines a local diffeomorphism and determines the Hamiltonian
$H \in C^2(T^*Q \times J)$ by
\[
H(q,p,t) = p\cdot \dot q - L(q,\dot q,t),
\]
with $\dot q$ understood as the inverse Legendre image of $(q,p)$.
\end{definition}

\begin{figure}[ht]
\centering
\begin{tikzpicture}[node distance=3.4cm, scale=1.05, every node/.style={scale=1.05}]
  \node[ccbox, align=center, minimum width=2.9cm] (C) at (0,2.9) {\textbf{Characteristic}\\Hamilton--Jacobi flow};
  \node[ccbox, align=center, minimum width=2.9cm] (L) at (-3.4,0) {\textbf{Lagrangian}\\Euler--Lagrange equations};
  \node[ccbox, align=center, minimum width=2.9cm] (H) at (3.4,0) {\textbf{Hamiltonian}\\Hamilton equations};

  \draw[ccarrow] (L) -- node[below, cclabel] {Legendre transform} (H);
  \draw[ccarrow] (H) -- node[above, cclabel] {inverse transform} (L);
  \draw[ccarrow] (L) -- node[left, cclabel] {extremal action} (C);
  \draw[ccarrow] (H) -- node[right, cclabel] {Hamilton--Jacobi} (C);
\end{tikzpicture}
\caption{Dynamics-side compatibility of the Lagrangian, Hamiltonian, and characteristic presentations for a regular variational system.}
\label{fig:dynamics-triad}
\end{figure}

\begin{theorem}[Dynamics-Side Triad Equivalence]
\label{thm:dynamics-triad}
\lean{dynamics_triad}
Let $L$ be a regular variational system in the sense of
Definition~\ref{def:regular-variational-system}, let $H$ be its Hamiltonian,
and let $S \in C^2(Q \times J)$ satisfy the Hamilton--Jacobi equation
\[
\partial_t S(q,t) + H\bigl(q,\partial_q S(q,t),t\bigr)=0
\]
on a simply connected domain. Then the following descriptions determine the same
local trajectories:
\begin{enumerate}[label=(\roman*), leftmargin=2em, itemsep=0.2em]
\item stationary curves of the action $\int L(q,\dot q,t)\,dt$;
\item solutions of Hamilton's equations for $H$;
\item characteristic curves of the Hamilton--Jacobi equation, given by
      \[
      p=\partial_q S,
      \qquad
      \dot q = \partial_p H(q,p,t).
      \]
\end{enumerate}
\end{theorem}

(Proof in Appendix~\ref{sec:appendix-dynamics}.)

\begin{proposition}[Characteristic Projection of Hamilton--Jacobi Dynamics]
\label{prop:hj-projection}
\lean{hj_projection}
Let $S$ satisfy the hypotheses of
Theorem~\ref{thm:dynamics-triad}, and define the time-dependent vector field
\[
V_t(q) := \partial_p H\bigl(q,\partial_q S(q,t),t\bigr)
\]
on $Q$. Then the integral curves of $V_t$ are exactly the configuration-space
projections of the Hamiltonian trajectories lying in the Lagrangian graph
\[
\Gamma_t := \{(q,\partial_q S(q,t)) : q \in Q\} \subset T^*Q.
\]
Consequently, the observer-side characteristic curves generated by
$\mathcal D_\chi = \partial_t + \mathcal L_{V_t}$ are precisely the projected
dynamics-side characteristics of the Hamilton--Jacobi equation.
\end{proposition}

(Proof in Appendix~\ref{sec:appendix-dynamics}.)

\begin{proposition}[Selection Principle for Realized Classes]
\label{prop:selection-problem}
\lean{selection_problem}
Let $\mathfrak R$ be any class of characteristic horizon realizations that also
arises from regular variational systems satisfying
Theorem~\ref{thm:dynamics-triad}. Then the existence of the Part~I calculus is
uniform on $\mathfrak R$. Any theorem asserting a unique admissible horizon, a
distinguished completion, or a canonical dimensionless invariant on
$\mathfrak R$ must therefore enter through additional \emph{selection
hypotheses} on the realized class rather than through realization alone.
\end{proposition}

(Proof in Appendix~\ref{sec:appendix-dynamics}.)

\subsection{From Realization to Forcing}

Proposition~\ref{prop:selection-problem} identifies the point at which
uniqueness may enter. The next definitions isolate the corresponding selection
data and admissible realized classes. The resulting theorem is constructive:
once a realized class has been specified, the remaining forcing problem is an
explicit admissibility problem on reduced operator data.

\begin{definition}[Selection datum at a fixed horizon]
\label{def:selection-datum}
\lean{SelectionDatum}
Fix a horizon $h$ and let
\[
R=(\cH^\bullet,d,\{\CCPi{k}\}_{k\in\bbN},\{U_t\}_{t\in J})
\]
be a characteristic horizon realization. A \emph{selection datum} for $R$ at
horizon $h$ consists of:
\begin{enumerate}[label=(\roman*), leftmargin=2em, itemsep=0.2em]
\item the observed sector $\CCHobs{h}(R):=\CCPi{h}\cH^\bullet$ together with the
      realized observed differential $\CCObs{d}{h}$ and defect
      $\CCDefect{d}{h}$;
\item a self-adjoint lower-bound operator $D_h$ on $\CCHobs{h}(R)$ derived from
      defect propagation and, when present, elimination/residual terms on the
      observed sector;
\item a finite-dimensional family $\mathcal Q_h(R)$ of candidate self-adjoint
      reduced laws or completion data on $\CCHobs{h}(R)$;
\item a unitary action of a finite group $G_h$ on $\CCHobs{h}(R)$ preserving
      $D_h$ and $\mathcal Q_h(R)$;
\item a family $\mathcal N_h$ of normalization, compatibility, or scale-fixing
      conditions imposed on $\mathcal Q_h(R)$.
\end{enumerate}
Such a datum will be denoted by $\mathsf S_h(R)$.
Locality and horizon compatibility are encoded in the choice of the
finite-dimensional candidate family $\mathcal Q_h(R)$ inside the realized
class.
\end{definition}

\begin{definition}[Horizon-preserving equivalence]
\label{def:horizon-preserving-equivalence}
\lean{HorizonPreservingEquivalence}
Two selection data
\[
\mathsf S_h(R)
\qquad\text{and}\qquad
\mathsf S_h(R')
\]
at the same horizon are \emph{horizon-preserving equivalent} if there exists a
unitary isomorphism
\[
W:\CCHobs{h}(R)\to \CCHobs{h}(R')
\]
such that:
\begin{enumerate}[label=(\roman*), leftmargin=2em, itemsep=0.2em]
\item $W\,\CCObs{d}{h} = \CCObs{d'}{h}\,W$ and
      $W\,\CCDefect{d}{h} = \CCDefect{d'}{h}\,W$;
\item $W D_h = D_h' W$;
\item $W\,\mathcal Q_h(R)\,W^{-1} = \mathcal Q_h(R')$;
\item $W$ intertwines the $G_h$-actions and preserves the normalization data
      $\mathcal N_h$.
\end{enumerate}
Equivalent selection data represent the same reduced observed law in different
coordinates.
\end{definition}

\begin{definition}[Admissible realized class]
\label{def:admissible-realized-class}
\lean{AdmissibleRealizedClass}
Fix a horizon $h$. An \emph{admissible realized class} at horizon $h$ is a
subclass $\mathfrak R_{\mathrm{adm}}(h)\subseteq \mathfrak R$ consisting of
those characteristic horizon realizations that are equipped with a selection
datum satisfying the following requirements:
\begin{enumerate}[label=(\roman*), leftmargin=2em, itemsep=0.2em]
\item the candidate family $\mathcal Q_h(R)$ is finite-dimensional;
\item the admissibility constraint of
      Definition~\ref{def:admissibility} is well posed for the derived
      lower-bound operator $D_h$ on $\CCHobs{h}(R)$;
\item $D_h$ and $\mathcal Q_h(R)$ are $G_h$-equivariant;
\item the normalization data $\mathcal N_h$ are invariant under
      horizon-preserving equivalence.
\end{enumerate}
If the horizon itself is part of the selection problem, one may further attach
a design objective on an appropriate Grassmannian as in
Definition~\ref{def:design-objective}.
\end{definition}

\begin{theorem}[Constructive Selection Interface for Realized Classes]
\label{thm:selection-interface}
\lean{selection_interface}
Fix a horizon $h$ and an admissible realized class
$\mathfrak R_{\mathrm{adm}}(h)$. Then for every
$R\in\mathfrak R_{\mathrm{adm}}(h)$ the forcing problem on the realized class
reduces to a finite and checkable selection problem:
\begin{enumerate}[label=(\roman*), leftmargin=2em, itemsep=0.2em]
\item the structural calculus of Part~I applies to $R$ on
      $\CCHobs{h}(R)$;
\item after symmetry reduction, the admissible reduced operator data are
      parameterized by finitely many block variables on the multiplicity spaces
      of the $G_h$-isotypic decomposition;
\item the completion-and-selection algebra of Chapter~5 returns exactly one of the following outputs:
      a unique admissible reduced datum, a minimal residual face
      certifying non-uniqueness, or an inconsistency certificate;
\item if horizon selection is included through a Grassmannian design objective,
      then optimal horizons exist in finite dimensions and satisfy the
      stationary commutator equation
      \[
      [\nabla \CCDesign{P^\star},P^\star]=0.
      \]
\end{enumerate}
In particular, realization plus admissibility converts the forcing question
into an explicit reduced problem on finite-dimensional operator data.
\end{theorem}

(Proof in Appendix~\ref{sec:appendix-dynamics}.)

\begin{definition}[Selected observed law]
\label{def:selected-observed-law}
\lean{selected_observed_law}
Let $R\in\mathfrak R_{\mathrm{adm}}(h)$. If the admissibility output in
Theorem~\ref{thm:selection-interface}(iii) is unique, and if any horizon
optimization included in the datum also yields a unique horizon up to
horizon-preserving equivalence, then the resulting reduced operator datum is
called the \emph{selected observed law} of $R$ at horizon $h$.
\end{definition}

\begin{corollary}[Selection-to-Forcing Principle]
\label{cor:selection-to-forcing}
\lean{selection_to_forcing}
Fix a horizon $h$ and an admissible realized class
$\mathfrak R_{\mathrm{adm}}(h)$. Suppose the admissible reduced selection data
modulo horizon-preserving equivalence consist of a single equivalence class.
Then the selected observed law is forced on
$\mathfrak R_{\mathrm{adm}}(h)$: every realization in
$\mathfrak R_{\mathrm{adm}}(h)$ yields the same selected observed law up to
horizon-preserving equivalence.
\end{corollary}

\begin{proof}
By Theorem~\ref{thm:selection-interface}, each realization in
$\mathfrak R_{\mathrm{adm}}(h)$ reduces to a finite admissibility problem on
its reduced operator data. By hypothesis, the admissible reduced data modulo
horizon-preserving equivalence form a singleton. Therefore the unique
selection output of Definition~\ref{def:selected-observed-law} is the same for
every realization up to that equivalence. Hence the observed law is forced on
the class, modulo change of coordinates preserving horizons, symmetry, and
normalization.
\end{proof}

\begin{definition}[Intrinsically equivariant realized class]\lean{IntrinsicallyEquivariantRealizedClass}
\label{def:intrinsically-equivariant-realized-class}
Fix a horizon $h$. An admissible realized class
$\mathfrak R_{\mathrm{adm}}(h)$ is \emph{intrinsically equivariant on the
canonical carrier} if each realization
$R\in\mathfrak R_{\mathrm{adm}}(h)$ is equipped with a unitary identification
\[
W_R:\CCHobs{h}(R)\to V_{\mathrm{slot}}
\]
onto the canonical six-slot carrier of
Definition~\ref{def:local-slot-carrier} such that:
\begin{enumerate}[label=(\roman*), leftmargin=2em, itemsep=0.2em]
\item the transported symmetry group $G_h$ is the intrinsic relabeling action
      of $S_4$ on $V_{\mathrm{slot}}$ from
      Definition~\ref{def:intrinsic-relabeling-action};
\item the transported lower-bound operator and candidate family are
      $S_4$-equivariant on $V_{\mathrm{slot}}$;
\item the normalization data are invariant under the same $S_4$ action.
\end{enumerate}
\end{definition}

\begin{theorem}[Intrinsic-equivariant forcing on the canonical carrier]\lean{intrinsic_equivariant_forcing}
\label{thm:intrinsic-equivariant-forcing}
Fix a horizon $h$ and let $\mathfrak R_{\mathrm{adm}}(h)$ be an intrinsically
equivariant realized class on the canonical carrier. Suppose the admissible
reduced selection data modulo horizon-preserving equivalence consist of a
single equivalence class. Then there exist unique scalars $a,b,c\in\bbR$ such
that for every $R\in\mathfrak R_{\mathrm{adm}}(h)$, the selected observed law
$L_R$ of $R$ satisfies
\[
W_R L_R W_R^{-1} e_\alpha
=
a\,e_\alpha
+
b\sum_{\beta\sim\alpha}e_\beta
+
c\sum_{\beta\perp\alpha}e_\beta
\qquad(\alpha\in\mathcal S_4).
\]
Equivalently, after transport to the canonical carrier, the selected observed
law is the same intrinsic-equivariant three-parameter operator for every
realization in the class. Its channel eigenvalues on the forced $1+2+3$
decomposition are
\[
\lambda_1=a+4b+c,\qquad
\lambda_2=a-2b+c,\qquad
\lambda_3=a-c.
\]
\end{theorem}

(Proof in Appendix~\ref{sec:appendix-dynamics}.)

\begin{definition}[Subgroup-broken realized class]\lean{SubgroupBrokenRealizedClass}
\label{def:subgroup-broken-realized-class}
Let $H\le S_4$ be a subgroup. An admissible realized class
$\mathfrak R_{\mathrm{adm}}(h)$ is \emph{$H$-equivariant on the canonical
carrier} if each realization admits a unitary identification
\[
W_R:\CCHobs{h}(R)\to V_{\mathrm{slot}}
\]
for which the transported symmetry, lower-bound operator, candidate family, and
normalization data are all invariant under the restricted action of $H$ on
$V_{\mathrm{slot}}$.
\end{definition}

\begin{proposition}[Finite subgroup-broken selection]\lean{finite_subgroup_broken_selection}
\label{prop:finite-subgroup-broken-selection}
Let $H\le S_4$ and let $\mathfrak R_{\mathrm{adm}}(h)$ be an
$H$-equivariant realized class on the canonical carrier. Then any selected
observed law on $\mathfrak R_{\mathrm{adm}}(h)$ decomposes blockwise on the
$H$-isotypic decomposition of $V_{\mathrm{slot}}$. In particular:
\begin{enumerate}[label=(\roman*), leftmargin=2em, itemsep=0.2em]
\item the corresponding reduced channel problem is still finite-dimensional;
\item the number of irreducible local channels is at most $6$;
\item the full intrinsic-equivariant case $H=S_4$ recovers the forced
      three-parameter $1+2+3$ law of
      Theorem~\ref{thm:intrinsic-equivariant-forcing}.
\end{enumerate}
\end{proposition}

\begin{proof}
After transport by $W_R$, the selected law is $H$-equivariant on the
six-dimensional carrier $V_{\mathrm{slot}}$. Theorem~\ref{thm:normal-form} and
Corollary~\ref{cor:isotypic-block} therefore reduce the law to block variables
on the multiplicity spaces of the $H$-isotypic decomposition. Because
$V_{\mathrm{slot}}$ has dimension $6$, there can be at most six nonzero
irreducible channels. This proves (i) and (ii). Statement (iii) is exactly the
special case $H=S_4$, where
Theorem~\ref{thm:intrinsic-equivariant-forcing} applies.
\end{proof}

\begin{definition}[Complex-type realized class]\lean{ComplexTypeRealizedClass}
\label{def:complex-type-realized-class}
Fix a horizon $h$. An admissible realized class
$\mathfrak R_{\mathrm{adm}}(h)$ is \emph{complex-type with scalar multiplicity
commutant} if, for every realization $R\in\mathfrak R_{\mathrm{adm}}(h)$:
\begin{enumerate}[label=(\roman*), leftmargin=2em, itemsep=0.2em]
\item the observed sector $\CCHobs{h}(R)$ carries a finite unitary symmetry
      action and an isotypic decomposition on the relevant selected-law sector;
\item the transported selected observed law and admissible selection data are
      equivariant for that action;
\item the selected-law sector carries a complex-type multiplicity package in
      the sense of Definition~\ref{def:complex-type-multiplicity-package};
\item on each isotypic channel, the induced multiplicity-space action of the
      selected-law family has scalar multiplicity commutant in the sense of
      Definition~\ref{def:scalar-multiplicity-commutant}.
\end{enumerate}
\end{definition}

\begin{corollary}[Realized phase-carrier class forcing]\lean{realized_phase_class_forcing}
\label{cor:realized-phase-class-forcing}
Fix a horizon $h$ and let $\mathfrak R_{\mathrm{adm}}(h)$ be a complex-type
realized class with scalar multiplicity commutant. Suppose the admissible
reduced selection data modulo horizon-preserving equivalence consist of a
single equivalence class. Then the selected law on
$\mathfrak R_{\mathrm{adm}}(h)$ carries a forced phase-carrier class modulo
phase-equivalence.
\end{corollary}

(Proof in Appendix~\ref{sec:appendix-dynamics}.)

\begin{definition}[Observable-irreducible realized class]\lean{ObservableIrreducibleRealizedClass}
\label{def:observable-irreducible-realized-class}
Fix a horizon $h$. An admissible realized class
$\mathfrak R_{\mathrm{adm}}(h)$ is \emph{observable-irreducible on the selected
channels} if, for every realization $R\in\mathfrak R_{\mathrm{adm}}(h)$ and
for every isotypic channel appearing in the selected-law sector:
\begin{enumerate}[label=(\roman*), leftmargin=2em, itemsep=0.2em]
\item the selected-law family acts on the corresponding multiplicity space
      through a family $\mathcal A_{\rho}(R)$ as in
      Definition~\ref{def:observable-irreducibility};
\item the family $\mathcal A_{\rho}(R)$ is observable-irreducible;
\item every selected self-adjoint closure or quadratic transport operator on
      that channel commutes with $\mathcal A_{\rho}(R)$.
\end{enumerate}
\end{definition}

\begin{corollary}[Realized self-adjoint scalarization on observable-irreducible channels]\lean{realized_self_adjoint_scalarization}
\label{cor:realized-self-adjoint-scalarization}
Fix a horizon $h$ and let $\mathfrak R_{\mathrm{adm}}(h)$ be an
observable-irreducible realized class. Then every selected self-adjoint
closure or quadratic transport operator on a selected channel is scalar on the
corresponding multiplicity space.
\end{corollary}

\begin{proof}
By Definition~\ref{def:observable-irreducible-realized-class}, the selected
self-adjoint operator on each channel commutes with an observable-irreducible
family on the corresponding multiplicity space. Corollary~\ref{cor:self-adjoint-commutant-scalarization}
then forces that operator to be a real scalar multiple of the identity on that
multiplicity space.
\end{proof}

\begin{corollary}[Realized isotropic symmetric-response scalarization]\lean{realized_isotropic_symmetric_response_scalarization}
\label{cor:realized-isotropic-symmetric-response-scalarization}
Fix a horizon $h$ and let $\mathfrak R_{\mathrm{adm}}(h)$ be an admissible
realized class. Assume that on a selected channel fiber the observable carrier
is the symmetric-tensor carrier $\mathrm{Sym}^2(F_h)$ of a real
$d$-dimensional inner-product space $F_h$ with $d\ge 2$, and that the selected
leading response operator
\[
\mathcal L_h:\mathrm{Sym}^2(F_h)\to \mathrm{Sym}^2(F_h)
\]
is equivariant for the induced orthogonal conjugation action of $O(F_h)$.
Then there are unique scalars $\alpha_h,\beta_h\in\bbR$ such that
\[
\mathcal L_h(S)
=
\alpha_h\Bigl(S-\frac{\operatorname{tr}(S)}{d}I\Bigr)
+
\beta_h\,\frac{\operatorname{tr}(S)}{d}I
\qquad\text{for every }S\in \mathrm{Sym}^2(F_h).
\]
In particular, the trace line $\bbR I$ and the traceless carrier
$\mathrm{Sym}^2_0(F_h)$ are the intrinsic scalar channels of the selected
response.
\end{corollary}

\begin{proof}
Choose an orthonormal identification $F_h\cong \bbR^d$. Under that
identification, the induced orthogonal conjugation action on
$\mathrm{Sym}^2(F_h)$ becomes the standard $O(d)$-action on
$\mathrm{Sym}^2(\bbR^d)$, and $\mathcal L_h$ satisfies the hypotheses of
Theorem~\ref{thm:orthogonal-symmetric-tensor-scalarization}. The resulting
scalars are independent of the chosen orthonormal identification because the
formula in that theorem is itself orthogonally invariant.
\end{proof}

\begin{corollary}[Realized phase forcing from real-type observable irreducibility]\lean{realized_phase_from_real_type}
\label{cor:realized-phase-from-real-type}
Fix a horizon $h$ and let $\mathfrak R_{\mathrm{adm}}(h)$ be an admissible
realized class such that:
\begin{enumerate}[label=(\roman*), leftmargin=2em, itemsep=0.2em]
\item on each relevant selected-law channel, the observed sector carries the
      data of a complex-type multiplicity package in the sense of
      Definition~\ref{def:complex-type-multiplicity-package};
\item on each relevant multiplicity space, the selected-law family is
      observable-irreducible and of real type in the sense of
      Definitions~\ref{def:observable-irreducibility} and
      \ref{def:real-type-observable-channel};
\item the admissible reduced selection data modulo horizon-preserving
      equivalence consist of a single equivalence class.
\end{enumerate}
Then the selected law on $\mathfrak R_{\mathrm{adm}}(h)$ carries a forced
phase-carrier class modulo phase-equivalence.
\end{corollary}

\begin{proof}
By Corollary~\ref{cor:selection-to-forcing}, the selected observed law is
forced on $\mathfrak R_{\mathrm{adm}}(h)$ up to horizon-preserving
equivalence. On each channel, the real-type observable irreducibility
hypothesis lets
Corollary~\ref{cor:real-type-irreducibility-collapses-phase-gauge} collapse
the residual phase gauge to a unique class modulo phase-equivalence. These
channelwise phase classes therefore agree across the realized class up to the
same equivalence.
\end{proof}

\begin{corollary}[Realized phase forcing from adjoint-closed real type]\lean{realized_phase_from_adjoint_closed_real_type}
\label{cor:realized-phase-from-adjoint-closed-real-type}
Fix a horizon $h$ and let $\mathfrak R_{\mathrm{adm}}(h)$ be an admissible
realized class such that:
\begin{enumerate}[label=(\roman*), leftmargin=2em, itemsep=0.2em]
\item on each relevant selected-law channel, the observed sector carries the
      data of a complex-type multiplicity package in the sense of
      Definition~\ref{def:complex-type-multiplicity-package};
\item on each relevant multiplicity space, the selected-law family is
      observable-irreducible, of real type, and adjoint-closed in generation
      in the sense of Definitions~\ref{def:observable-irreducibility},
      \ref{def:real-type-observable-channel}, and
      \ref{def:adjoint-closed-generated-observable-algebra};
\item the admissible reduced selection data modulo horizon-preserving
      equivalence consist of a single equivalence class.
\end{enumerate}
Then the selected law on $\mathfrak R_{\mathrm{adm}}(h)$ carries a forced
phase-carrier class modulo phase-equivalence.
\end{corollary}

\begin{proof}
By Corollary~\ref{cor:selection-to-forcing}, the selected observed law is
forced on $\mathfrak R_{\mathrm{adm}}(h)$ up to horizon-preserving
equivalence. On each channel, the adjoint-closed real-type hypothesis lets
Corollary~\ref{cor:adjoint-closed-real-type-phase} collapse the residual phase
gauge to a unique class modulo phase-equivalence. These channelwise phase
classes therefore agree across the realized class up to the same equivalence.
\end{proof}

\begin{corollary}[Realized phase forcing from a rank-one observable seed]\lean{realized_phase_from_rank_one_seed}
\label{cor:realized-phase-from-rank-one-seed}
Fix a horizon $h$ and let $\mathfrak R_{\mathrm{adm}}(h)$ be an admissible
realized class such that:
\begin{enumerate}[label=(\roman*), leftmargin=2em, itemsep=0.2em]
\item on each relevant selected-law channel, the observed sector carries the
      data of a complex-type multiplicity package in the sense of
      Definition~\ref{def:complex-type-multiplicity-package};
\item on each relevant multiplicity space, the selected-law family is
      observable-irreducible, adjoint-closed in generation, and carries a
      rank-one observable seed in the sense of
      Definitions~\ref{def:observable-irreducibility},
      \ref{def:adjoint-closed-generated-observable-algebra}, and
      \ref{def:rank-one-observable-seed};
\item the admissible reduced selection data modulo horizon-preserving
      equivalence consist of a single equivalence class.
\end{enumerate}
Then the selected law on $\mathfrak R_{\mathrm{adm}}(h)$ carries a forced
phase-carrier class modulo phase-equivalence.
\end{corollary}

\begin{proof}
By Corollary~\ref{cor:selection-to-forcing}, the selected observed law is
forced on $\mathfrak R_{\mathrm{adm}}(h)$ up to horizon-preserving
equivalence. On each channel, the rank-one seed hypothesis lets
Corollary~\ref{cor:rank-one-seed-phase} collapse the residual phase gauge to a
unique class modulo phase-equivalence. These channelwise phase classes
therefore agree across the realized class up to the same equivalence.
\end{proof}

\begin{corollary}[Realized phase forcing from simple spectrum]\lean{realized_phase_from_simple_spectrum}
\label{cor:realized-phase-from-simple-spectrum}
Fix a horizon $h$ and let $\mathfrak R_{\mathrm{adm}}(h)$ be an admissible
realized class such that:
\begin{enumerate}[label=(\roman*), leftmargin=2em, itemsep=0.2em]
\item on each relevant selected-law channel, the observed sector carries the
      data of a complex-type multiplicity package in the sense of
      Definition~\ref{def:complex-type-multiplicity-package};
\item on each relevant multiplicity space, the selected-law family is
      observable-irreducible, adjoint-closed in generation, and the generated
      algebra contains a self-adjoint simple-spectrum observable in the sense
      of Definition~\ref{def:simple-spectrum-observable};
\item the admissible reduced selection data modulo horizon-preserving
      equivalence consist of a single equivalence class.
\end{enumerate}
Then the selected law on $\mathfrak R_{\mathrm{adm}}(h)$ carries a forced
phase-carrier class modulo phase-equivalence.
\end{corollary}

\begin{proof}
By Corollary~\ref{cor:selection-to-forcing}, the selected observed law is
forced on $\mathfrak R_{\mathrm{adm}}(h)$ up to horizon-preserving
equivalence. On each channel, the simple-spectrum hypothesis lets
Corollary~\ref{cor:simple-spectrum-phase} collapse the residual phase gauge to
a unique class modulo phase-equivalence. These channelwise phase classes
therefore agree across the realized class up to the same equivalence.
\end{proof}

\begin{corollary}[Realized phase forcing from observable completeness]\lean{realized_phase_from_observable_completeness}
\label{cor:realized-phase-from-observable-completeness}
Fix a horizon $h$ and let $\mathfrak R_{\mathrm{adm}}(h)$ be an admissible
realized class such that:
\begin{enumerate}[label=(\roman*), leftmargin=2em, itemsep=0.2em]
\item on each relevant selected-law channel, the observed sector carries the
      data of a complex-type multiplicity package in the sense of
      Definition~\ref{def:complex-type-multiplicity-package};
\item on each relevant multiplicity space, the selected-law family is
      observable-complete in the sense of
      Definition~\ref{def:observable-completeness};
\item the admissible reduced selection data modulo horizon-preserving
      equivalence consist of a single equivalence class.
\end{enumerate}
Then the selected law on $\mathfrak R_{\mathrm{adm}}(h)$ carries a forced
phase-carrier class modulo phase-equivalence.
\end{corollary}

\begin{proof}
By Corollary~\ref{cor:selection-to-forcing}, the selected observed law is
forced on $\mathfrak R_{\mathrm{adm}}(h)$ up to horizon-preserving
equivalence. On each channel, the observable-completeness hypothesis lets
Corollary~\ref{cor:observable-completeness-collapses-phase-gauge} collapse the
phase gauge to a unique class modulo phase-equivalence. These channelwise phase
classes therefore agree across the realized class up to the same equivalence.
\end{proof}

\begin{corollary}[Realized phase forcing from a matrix-unit scaffold]\lean{realized_phase_from_matrix_units}
\label{cor:realized-phase-from-matrix-units}
Fix a horizon $h$ and let $\mathfrak R_{\mathrm{adm}}(h)$ be an admissible
realized class such that:
\begin{enumerate}[label=(\roman*), leftmargin=2em, itemsep=0.2em]
\item on each relevant selected-law channel, the observed sector carries the
      data of a complex-type multiplicity package in the sense of
      Definition~\ref{def:complex-type-multiplicity-package};
\item on each relevant multiplicity space, the selected-law family contains a
      matrix-unit observable scaffold in the sense of
      Definition~\ref{def:matrix-unit-observable-scaffold};
\item the admissible reduced selection data modulo horizon-preserving
      equivalence consist of a single equivalence class.
\end{enumerate}
Then the selected law on $\mathfrak R_{\mathrm{adm}}(h)$ carries a forced
phase-carrier class modulo phase-equivalence.
\end{corollary}

\begin{proof}
By Corollary~\ref{cor:selection-to-forcing}, the selected observed law is
forced on $\mathfrak R_{\mathrm{adm}}(h)$ up to horizon-preserving
equivalence. On each channel, the matrix-unit scaffold hypothesis lets
Corollary~\ref{cor:matrix-unit-scaffold-phase} collapse the phase gauge to a
unique class modulo phase-equivalence. These channelwise phase classes
therefore agree across the realized class up to the same equivalence.
\end{proof}

\begin{corollary}[Realized finite observable rigidity and phase forcing]\lean{realized_finite_observable_rigidity}
\label{cor:realized-finite-observable-rigidity}
Fix a horizon $h$ and let $\mathfrak R_{\mathrm{adm}}(h)$ be an admissible
realized class such that:
\begin{enumerate}[label=(\roman*), leftmargin=2em, itemsep=0.2em]
\item on each relevant selected-law channel, the observed sector carries the
      data of a complex-type multiplicity package in the sense of
      Definition~\ref{def:complex-type-multiplicity-package};
\item on each relevant multiplicity space, the selected-law family has
      generated algebra that is adjoint-closed and observable-irreducible;
\item on each relevant multiplicity space, at least one of the equivalent
      conditions of Theorem~\ref{thm:finite-observable-rigidity} holds;
\item the admissible reduced selection data modulo horizon-preserving
      equivalence consist of a single equivalence class.
\end{enumerate}
Then on every relevant selected-law channel all four conditions of
Theorem~\ref{thm:finite-observable-rigidity} hold, and the selected law on
$\mathfrak R_{\mathrm{adm}}(h)$ carries a forced phase-carrier class modulo
phase-equivalence.
\end{corollary}

(Proof in Appendix~\ref{sec:appendix-dynamics}.)

\begin{definition}[Phase-generated realized class]\lean{PhaseGeneratedRealizedClass}
\label{def:phase-generated-realized-class}
Fix a horizon $h$. An admissible realized class
$\mathfrak R_{\mathrm{adm}}(h)$ is \emph{phase-generated on the selected
channels} if:
\begin{enumerate}[label=(\roman*), leftmargin=2em, itemsep=0.2em]
\item on each relevant selected-law channel, the observed sector carries the
      data of a complex-type multiplicity package in the sense of
      Definition~\ref{def:complex-type-multiplicity-package};
\item on each relevant multiplicity space, the selected-law family is
      phase-generating in the sense of
      Definition~\ref{def:phase-generating-observable-family};
\item the admissible reduced selection data modulo horizon-preserving
      equivalence consist of a single equivalence class.
\end{enumerate}
\end{definition}

\begin{corollary}[Realized phase generation forcing]\lean{realized_phase_generation_forcing}
\label{cor:realized-phase-generation-forcing}
Fix a horizon $h$ and let $\mathfrak R_{\mathrm{adm}}(h)$ be a phase-generated
realized class. Then the selected law on $\mathfrak R_{\mathrm{adm}}(h)$
carries a forced phase-carrier class modulo phase-equivalence.
\end{corollary}

\begin{proof}
By Corollary~\ref{cor:selection-to-forcing}, the selected observed law is
forced on $\mathfrak R_{\mathrm{adm}}(h)$ up to horizon-preserving
equivalence. On each relevant channel,
Theorem~\ref{thm:phase-generation-forces-phase-class} collapses the residual
phase gauge to a unique class modulo phase-equivalence. These channelwise
phase classes therefore agree across the realized class up to the same
equivalence.
\end{proof}

\begin{corollary}[Realized phase generation implies phase rigidity]\lean{realized_phase_generation_implies_rigidity}
\label{cor:realized-phase-generation-implies-rigidity}
Fix a horizon $h$ and let $\mathfrak R_{\mathrm{adm}}(h)$ be a phase-generated
realized class. Then $\mathfrak R_{\mathrm{adm}}(h)$ is phase-rigid on the
selected channels in the sense of
Definition~\ref{def:phase-rigid-realized-class}.
\end{corollary}

(Proof in Appendix~\ref{sec:appendix-dynamics}.)

\begin{definition}[Phase-rigid realized class]\lean{PhaseRigidRealizedClass}
\label{def:phase-rigid-realized-class}
Fix a horizon $h$. An admissible realized class
$\mathfrak R_{\mathrm{adm}}(h)$ is \emph{phase-rigid on the selected channels}
if:
\begin{enumerate}[label=(\roman*), leftmargin=2em, itemsep=0.2em]
\item on each relevant selected-law channel, the observed sector carries the
      data of a complex-type multiplicity package in the sense of
      Definition~\ref{def:complex-type-multiplicity-package};
\item on each relevant multiplicity space, the selected-law family is
      phase-rigid in the sense of
      Definition~\ref{def:phase-rigid-observable-family};
\item the admissible reduced selection data modulo horizon-preserving
      equivalence consist of a single equivalence class.
\end{enumerate}
\end{definition}

\begin{corollary}[Realized phase rigidity forcing]\lean{realized_phase_rigidity_forcing}
\label{cor:realized-phase-rigidity-forcing}
Fix a horizon $h$ and let $\mathfrak R_{\mathrm{adm}}(h)$ be a phase-rigid
realized class. Then the selected law on $\mathfrak R_{\mathrm{adm}}(h)$
carries a forced phase-carrier class modulo phase-equivalence.
\end{corollary}

\begin{proof}
By Corollary~\ref{cor:selection-to-forcing}, the selected observed law is
forced on $\mathfrak R_{\mathrm{adm}}(h)$ up to horizon-preserving
equivalence. On each channel,
Corollary~\ref{cor:phase-rigidity-collapses-phase-gauge} collapses the
residual phase gauge to a unique class modulo phase-equivalence. These
channelwise phase classes therefore agree across the realized class up to the
same equivalence.
\end{proof}

\begin{definition}[Observable-rigid realized class]\lean{ObservableRigidRealizedClass}
\label{def:observable-rigid-realized-class}
Fix a horizon $h$. An admissible realized class
$\mathfrak R_{\mathrm{adm}}(h)$ is \emph{observable-rigid on the selected
channels} if:
\begin{enumerate}[label=(\roman*), leftmargin=2em, itemsep=0.2em]
\item $\mathfrak R_{\mathrm{adm}}(h)$ is phase-rigid on the selected channels
      in the sense of Definition~\ref{def:phase-rigid-realized-class};
\item $\mathfrak R_{\mathrm{adm}}(h)$ is observable-irreducible on the
      selected channels in the sense of
      Definition~\ref{def:observable-irreducible-realized-class}.
\end{enumerate}
\end{definition}

\begin{corollary}[Observable-complete classes are observable-rigid]\lean{observable_complete_implies_observable_rigid}
\label{cor:observable-complete-implies-observable-rigid}
Fix a horizon $h$ and let $\mathfrak R_{\mathrm{adm}}(h)$ be an admissible
realized class such that:
\begin{enumerate}[label=(\roman*), leftmargin=2em, itemsep=0.2em]
\item on each relevant selected-law channel, the observed sector carries the
      data of a complex-type multiplicity package in the sense of
      Definition~\ref{def:complex-type-multiplicity-package};
\item on each relevant multiplicity space, the selected-law family is
      observable-complete in the sense of
      Definition~\ref{def:observable-completeness};
\item the admissible reduced selection data modulo horizon-preserving
      equivalence consist of a single equivalence class.
\end{enumerate}
Then $\mathfrak R_{\mathrm{adm}}(h)$ is observable-rigid on the selected
channels in the sense of
Definition~\ref{def:observable-rigid-realized-class}.
\end{corollary}

(Proof in Appendix~\ref{sec:appendix-dynamics}.)

\begin{corollary}[Finite-observable-rigidity classes are observable-rigid]\lean{finite_observable_rigidity_implies_observable_rigid}
\label{cor:finite-observable-rigidity-implies-observable-rigid}
Fix a horizon $h$ and let $\mathfrak R_{\mathrm{adm}}(h)$ be an admissible
realized class such that:
\begin{enumerate}[label=(\roman*), leftmargin=2em, itemsep=0.2em]
\item on each relevant selected-law channel, the observed sector carries the
      data of a complex-type multiplicity package in the sense of
      Definition~\ref{def:complex-type-multiplicity-package};
\item on each relevant multiplicity space, the selected-law family has
      generated algebra that is adjoint-closed and observable-irreducible;
\item on each relevant multiplicity space, at least one of the equivalent
      conditions of Theorem~\ref{thm:finite-observable-rigidity} holds;
\item the admissible reduced selection data modulo horizon-preserving
      equivalence consist of a single equivalence class.
\end{enumerate}
Then $\mathfrak R_{\mathrm{adm}}(h)$ is observable-rigid on the selected
channels in the sense of
Definition~\ref{def:observable-rigid-realized-class}.
\end{corollary}

\begin{proof}
On each relevant multiplicity space,
Theorem~\ref{thm:finite-observable-rigidity} upgrades item (iii) to
observable completeness. Therefore
Corollary~\ref{cor:observable-complete-implies-observable-rigid} applies.
\end{proof}

\begin{definition}[Observable-orbit generated realized class]\lean{ObservableOrbitGeneratedRealizedClass}
\label{def:observable-orbit-generated-realized-class}
Fix a horizon $h$. An admissible realized class
$\mathfrak R_{\mathrm{adm}}(h)$ is \emph{observable-orbit generated on the
selected channels} if:
\begin{enumerate}[label=(\roman*), leftmargin=2em, itemsep=0.2em]
\item on each relevant selected-law channel, the observed sector carries the
      data of a complex-type multiplicity package in the sense of
      Definition~\ref{def:complex-type-multiplicity-package};
\item on each relevant multiplicity space $M_{\rho}$, there is an orthogonal
      action
      \[
      V_{\rho}:K_{\rho}\to O(M_{\rho})
      \]
      that is transitive on the unit sphere of $M_{\rho}$;
\item on each relevant multiplicity space $M_{\rho}$, there is a rank-one
      orthogonal projection
      \[
      P_{0,\rho}:M_{\rho}\to M_{\rho}
      \]
      such that the selected-law family contains every conjugate
      \[
      V_{\rho}(g)P_{0,\rho}V_{\rho}(g)^*
      \qquad(g\in K_{\rho});
      \]
\item the admissible reduced selection data modulo horizon-preserving
      equivalence consist of a single equivalence class.
\end{enumerate}
\end{definition}

\begin{corollary}[Observable-orbit generated classes are observable-rigid]\lean{observable_orbit_generated_implies_observable_rigid}
\label{cor:observable-orbit-generated-implies-observable-rigid}
Fix a horizon $h$ and let $\mathfrak R_{\mathrm{adm}}(h)$ be an
observable-orbit generated realized class on the selected channels. Then
$\mathfrak R_{\mathrm{adm}}(h)$ is observable-rigid on the selected channels in
the sense of Definition~\ref{def:observable-rigid-realized-class}.
\end{corollary}

\begin{proof}
For each relevant selected-law channel, items (ii)--(iii) place the
selected-law family inside the scope of
Corollary~\ref{cor:sphere-transitive-orbit-observable-completeness}, so that
family is observable-complete on the corresponding multiplicity space.
Together with items (i) and (iv), this is exactly the hypothesis package of
Corollary~\ref{cor:observable-complete-implies-observable-rigid}. Therefore
$\mathfrak R_{\mathrm{adm}}(h)$ is observable-rigid on the selected channels.
\end{proof}

\begin{corollary}[Adjoint-closed real-type classes are observable-rigid]\lean{adjoint_closed_real_type_implies_observable_rigid}
\label{cor:adjoint-closed-real-type-implies-observable-rigid}
Fix a horizon $h$ and let $\mathfrak R_{\mathrm{adm}}(h)$ be an admissible
realized class such that:
\begin{enumerate}[label=(\roman*), leftmargin=2em, itemsep=0.2em]
\item on each relevant selected-law channel, the observed sector carries the
      data of a complex-type multiplicity package in the sense of
      Definition~\ref{def:complex-type-multiplicity-package};
\item on each relevant multiplicity space, the selected-law family is
      observable-irreducible, of real type, and adjoint-closed in generation
      in the sense of Definitions~\ref{def:observable-irreducibility},
      \ref{def:real-type-observable-channel}, and
      \ref{def:adjoint-closed-generated-observable-algebra};
\item the admissible reduced selection data modulo horizon-preserving
      equivalence consist of a single equivalence class.
\end{enumerate}
Then $\mathfrak R_{\mathrm{adm}}(h)$ is observable-rigid on the selected
channels in the sense of
Definition~\ref{def:observable-rigid-realized-class}.
\end{corollary}

\begin{proof}
For each relevant selected-law channel, item (ii) already gives
observable-irreducibility and phase rigidity in the sense of
Definition~\ref{def:phase-rigid-observable-family}. Together with items (i)
and (iii), this supplies the phase-rigid and observable-irreducible realized
class conditions. Therefore both clauses of
Definition~\ref{def:observable-rigid-realized-class} hold.
\end{proof}

\begin{corollary}[Scalar-commutant complex-type classes are observable-rigid]\lean{scalar_commutant_complex_type_implies_observable_rigid}
\label{cor:scalar-commutant-complex-type-implies-observable-rigid}
Fix a horizon $h$ and let $\mathfrak R_{\mathrm{adm}}(h)$ be an admissible
realized class such that:
\begin{enumerate}[label=(\roman*), leftmargin=2em, itemsep=0.2em]
\item $\mathfrak R_{\mathrm{adm}}(h)$ is complex-type with scalar multiplicity
      commutant in the sense of
      Definition~\ref{def:complex-type-realized-class};
\item $\mathfrak R_{\mathrm{adm}}(h)$ is observable-irreducible on the
      selected channels in the sense of
      Definition~\ref{def:observable-irreducible-realized-class};
\item the admissible reduced selection data modulo horizon-preserving
      equivalence consist of a single equivalence class.
\end{enumerate}
Then $\mathfrak R_{\mathrm{adm}}(h)$ is observable-rigid on the selected
channels in the sense of
Definition~\ref{def:observable-rigid-realized-class}.
\end{corollary}

(Proof in Appendix~\ref{sec:appendix-dynamics}.)

\begin{corollary}[Generated observable classes are observable-rigid]\lean{generated_observable_implies_rigid}
\label{cor:generated-observable-implies-rigid}
Fix a horizon $h$ and let $\mathfrak R_{\mathrm{adm}}(h)$ be an admissible
realized class that is phase-generated and observable-irreducible on the
selected channels. Then $\mathfrak R_{\mathrm{adm}}(h)$ is observable-rigid on
the selected channels in the sense of
Definition~\ref{def:observable-rigid-realized-class}.
\end{corollary}

\begin{proof}
By Corollary~\ref{cor:realized-phase-generation-implies-rigidity},
$\mathfrak R_{\mathrm{adm}}(h)$ is phase-rigid on the selected channels. The
observable-irreducible hypothesis is already part of the assumptions.
Therefore both clauses of
Definition~\ref{def:observable-rigid-realized-class} hold.
\end{proof}

\begin{definition}[Transport-order disciplined realized class]\lean{TransportOrderDisciplinedRealizedClass}
\label{def:transport-order-disciplined-realized-class}
Fix a horizon $h$. An admissible realized class
$\mathfrak R_{\mathrm{adm}}(h)$ is \emph{transport-order disciplined} if, for
every realization $R\in\mathfrak R_{\mathrm{adm}}(h)$, the admissible selected
transport bases at horizon $h$ form a transport-order filtered admissible
family in the sense of Definition~\ref{def:transport-order-filtered-family},
and the minimal-order admissible subclass is singleton modulo the chosen
realized equivalence relation.
\end{definition}

\begin{corollary}[Realized transport-order forcing]\lean{realized_transport_order_forcing}
\label{cor:realized-transport-order-forcing}
Fix a horizon $h$ and let $\mathfrak R_{\mathrm{adm}}(h)$ be a transport-order
disciplined realized class. Then the selected law on
$\mathfrak R_{\mathrm{adm}}(h)$ carries a forced transport-order class. If the
selected transport bases admit exact Schur transport jets, those jets satisfy
the canonical coefficient bounds of
Corollary~\ref{cor:minimal-transport-order-controls-jet}.
\end{corollary}

\begin{proof}
For each realization, Theorem~\ref{thm:minimal-admissible-transport-order}
forces the distinguished transport-order class because the minimal-order
admissible subclass is singleton modulo the chosen realized equivalence
relation. The jet bound is then exactly
Corollary~\ref{cor:minimal-transport-order-controls-jet}.
\end{proof}

\begin{definition}[Selected channel transport system]\lean{SelectedChannelTransportSystem}
\label{def:selected-channel-transport-system}
Fix a horizon $h$. A realization in an admissible realized class carries a
\emph{selected channel transport system} if:
\begin{enumerate}[label=(\roman*), leftmargin=2em, itemsep=0.2em]
\item there is a finite observer graph $\Lambda_h=(\Lambda_h^0,\Lambda_h^1)$;
\item to each site $x\in\Lambda_h^0$ is attached a selected channel fiber
      $F_{h,x}$ inside the observed sector;
\item to each edge $(x,y)\in\Lambda_h^1$ is attached a unitary transport map
      \[
      U_{xy}:F_{h,y}\to F_{h,x}
      \]
      compatible with the selected-law symmetry data.
\end{enumerate}
\end{definition}

\begin{corollary}[Realized local gauge covariance on scalar-commutant channels]\lean{selected_channel_gauge_covariance}
\label{cor:realized-local-gauge-covariance}
Fix a horizon $h$ and let $\mathfrak R_{\mathrm{adm}}(h)$ be a complex-type
realized class with scalar multiplicity commutant. Assume each realization
$R\in\mathfrak R_{\mathrm{adm}}(h)$ carries a selected channel transport system
in the sense of Definition~\ref{def:selected-channel-transport-system}. Then on
every such selected channel transport system:
\begin{enumerate}[label=(\roman*), leftmargin=2em, itemsep=0.2em]
\item the transport energy is invariant under independent local unitary frame
      changes on the selected channel fibers;
\item if the selected channel fibers carry preserved unit volume forms, the
      allowed local frame group reduces to the unimodular subgroup.
\end{enumerate}
\end{corollary}

\begin{proof}
By Corollary~\ref{cor:realized-phase-class-forcing}, the selected-law channels
carry forced phase-carrier classes modulo phase-equivalence, and by
Definition~\ref{def:complex-type-realized-class} the local selected-law data
have scalar multiplicity commutant on each channel. Therefore the selected
channel fibers fit the hypotheses of
Theorem~\ref{thm:local-unitary-covariance}, which proves (i). If preserved unit
volume forms are also part of the realized data, then
Corollary~\ref{cor:unimodular-channel-reduction} gives (ii).
\end{proof}

\begin{corollary}[Realized reversible recursion forcing on selected channels]\lean{selected_channel_recursive_forcing}
\label{cor:realized-reversible-recursion-forcing}
Fix a horizon $h$ and let $\mathfrak R_{\mathrm{adm}}(h)$ be a transport-order
disciplined realized class. Suppose that on a selected channel fiber $F_h$:
\begin{enumerate}[label=(\roman*), leftmargin=2em, itemsep=0.2em]
\item the admissible transport bases satisfy the hypotheses of
      Theorem~\ref{thm:order-one-forcing-criterion};
\item the forced one-step transport class is represented by a bounded transport
      base
      \[
      B_h:F_h\to F_h;
      \]
\end{enumerate}
Then:
\begin{enumerate}[label=(\alph*), leftmargin=2em, itemsep=0.2em]
\item the distinguished admissible transport-order class on $F_h$ is $1$;
\item every bi-infinite selected channel transport sequence
      $x=(x_n)_{n\in\bbZ}$ satisfies the forced recursion constraint
      \[
      \mathcal R_{B_h}x=0,
      \qquad\text{equivalently}\qquad
      x_{n+1}=B_hx_n-x_{n-1};
      \].
\end{enumerate}
\end{corollary}

\begin{proof}
Hypothesis (i) and Theorem~\ref{thm:order-one-forcing-criterion} give (a). Once
the forced order-one class is represented by the bounded transport base $B_h$,
Theorem~\ref{thm:reversible-transport-recursion} gives the recursion identity
in (b).
\end{proof}

\begin{corollary}[Realized phase-compatible recursion forcing]\lean{selected_channel_recursive_forcing}
\label{cor:realized-recursion-forcing}
Fix a horizon $h$ and let $\mathfrak R_{\mathrm{adm}}(h)$ be a transport-order
disciplined realized class. Suppose that on a selected channel fiber $F_h$:
\begin{enumerate}[label=(\roman*), leftmargin=2em, itemsep=0.2em]
\item the hypotheses of
      Corollary~\ref{cor:realized-reversible-recursion-forcing} hold;
\item the selected channel also carries a forced phase-carrier class.
\end{enumerate}
Then every bi-infinite selected channel transport sequence
\[
x=(x_n)_{n\in\bbZ}
\]
satisfies the forced recursion constraint
\[
\mathcal R_{B_h}x=0,
\qquad\text{equivalently}\qquad
x_{n+1}=B_hx_n-x_{n-1},
\]
and this recursion constraint is phase-compatible.
\end{corollary}

\begin{proof}
Corollary~\ref{cor:realized-reversible-recursion-forcing} gives the recursion
law. The forced phase-carrier class together with
Proposition~\ref{prop:transport-recursion-compatibility} yields the stated
phase compatibility.
\end{proof}

\begin{definition}[Isotropic selected channel system]\lean{IsotropicSelectedChannelSystem}
\label{def:isotropic-selected-channel-system}
Fix a horizon $h$. A selected channel transport system is \emph{isotropic} on a
channel fiber $F_h$ if:
\begin{enumerate}[label=(\roman*), leftmargin=2em, itemsep=0.2em]
\item there is an orthogonal action
      \[
      U_h:K_h\to O(F_h)
      \]
      of a symmetry group $K_h$;
\item the representation of $K_h$ on $F_h$ is irreducible as a real module;
\item the selected one-step transport base
      \[
      B_h:F_h\to F_h
      \]
      commutes with the action of $K_h$.
\end{enumerate}
\end{definition}

\begin{corollary}[Realized isotropic quadratic transport]\lean{selected_channel_scalarization}
\label{cor:realized-isotropic-quadratic-transport}
Fix a horizon $h$ and let $\mathfrak R_{\mathrm{adm}}(h)$ be a transport-order
disciplined realized class. Suppose a selected channel fiber $F_h$ carries an
isotropic selected channel system with forced order-one transport base $B_h$.
Then there exists a unique scalar
\[
\kappa_h\ge 0
\]
such that
\[
B_h^*B_h=\kappa_h\,I_{F_h}.
\]
Equivalently, the selected quadratic transport energy on that channel is
\[
E_{B_h}(v)=\kappa_h\,\|v\|^2
\qquad(v\in F_h).
\]
\end{corollary}

\begin{proof}
The selected channel satisfies the hypotheses of
Theorem~\ref{thm:isotropic-quadratic-transport}. Applying that theorem to the
forced transport base $B_h$ gives the stated scalarization.
\end{proof}

\begin{definition}[Sphere-transitive selected channel system]\lean{SphereTransitiveSelectedChannelSystem}
\label{def:sphere-transitive-selected-channel-system}
Fix a horizon $h$. A selected channel transport system is
\emph{sphere-transitive} on a channel fiber $F_h$ if:
\begin{enumerate}[label=(\roman*), leftmargin=2em, itemsep=0.2em]
\item there is an orthogonal action
      \[
      U_h:K_h\to O(F_h)
      \]
      of a symmetry group $K_h$;
\item for any unit vectors $u,v\in F_h$ there exists $g\in K_h$ such that
      \[
      U_h(g)u=v;
      \]
\item the selected one-step transport base
      \[
      B_h:F_h\to F_h
      \]
      commutes with the action of $K_h$.
\end{enumerate}
\end{definition}

\begin{corollary}[Realized sphere-transitive quadratic transport]\lean{selected_channel_scalarization}
\label{cor:realized-sphere-transitive-quadratic-transport}
Fix a horizon $h$ and let $\mathfrak R_{\mathrm{adm}}(h)$ be a transport-order
disciplined realized class. Suppose a selected channel fiber $F_h$ carries a
sphere-transitive selected channel system with forced order-one transport base
$B_h$. Then there exists a unique scalar
\[
\kappa_h\ge 0
\]
such that
\[
B_h^*B_h=\kappa_h\,I_{F_h}.
\]
\end{corollary}

\begin{proof}
The selected channel satisfies the hypotheses of
Definition~\ref{def:sphere-transitive-transport-channel}. Therefore
Corollary~\ref{cor:sphere-transitive-quadratic-transport} applies directly to
the forced transport base $B_h$.
\end{proof}

\begin{definition}[Full orthogonal selected channel system]\lean{FullOrthogonalSelectedChannelSystem}
\label{def:full-orthogonal-selected-channel-system}
Fix a horizon $h$. A selected channel transport system is
\emph{fully orthogonal} on a channel fiber $F_h$ if:
\begin{enumerate}[label=(\roman*), leftmargin=2em, itemsep=0.2em]
\item the symmetry action on $F_h$ is the tautological action of the full
      orthogonal group
      \[
      O(F_h)\to O(F_h);
      \]
\item the selected one-step transport base
      \[
      B_h:F_h\to F_h
      \]
      commutes with every orthogonal transformation of $F_h$.
\end{enumerate}
\end{definition}

\begin{corollary}[Realized full orthogonal quadratic transport]\lean{selected_channel_scalarization}
\label{cor:realized-full-orthogonal-quadratic-transport}
Fix a horizon $h$ and let $\mathfrak R_{\mathrm{adm}}(h)$ be a transport-order
disciplined realized class. Suppose a selected channel fiber $F_h$ carries a
full orthogonal selected channel system with forced order-one transport base
$B_h$. Then there exists a unique scalar
\[
\kappa_h\ge 0
\]
such that
\[
B_h^*B_h=\kappa_h\,I_{F_h}.
\]
\end{corollary}

\begin{proof}
The selected channel satisfies the hypotheses of
Definition~\ref{def:full-orthogonal-transport-channel}. Therefore
Corollary~\ref{cor:full-orthogonal-quadratic-transport} applies directly to
the forced transport base $B_h$.
\end{proof}

\begin{definition}[Metric-isotropic selected channel system]\lean{MetricIsotropicSelectedChannelSystem}
\label{def:metric-isotropic-selected-channel-system}
Fix a horizon $h$. A selected channel transport system is
\emph{metric-isotropic} on a channel fiber $F_h$ if the selected one-step
transport base
\[
B_h:F_h\to F_h
\]
satisfies
\[
\|B_hu\|^2=\|B_hv\|^2
\]
for every pair of unit vectors $u,v\in F_h$.
\end{definition}

\begin{corollary}[Realized metric-isotropic quadratic transport]\lean{selected_channel_scalarization}
\label{cor:realized-metric-isotropic-quadratic-transport}
Fix a horizon $h$ and let $\mathfrak R_{\mathrm{adm}}(h)$ be a transport-order
disciplined realized class. Suppose a selected channel fiber $F_h$ carries a
metric-isotropic selected channel system with forced order-one transport base
$B_h$. Then there exists a unique scalar
\[
\kappa_h\ge 0
\]
such that
\[
B_h^*B_h=\kappa_h\,I_{F_h}.
\]
\end{corollary}

\begin{proof}
The metric-isotropy hypothesis places $B_h$ inside the scope of
Corollary~\ref{cor:metric-isotropic-transport}, which gives the stated
scalarization.
\end{proof}

\begin{definition}[Direction-frame selected channel system]\lean{DirectionFrameSelectedChannelSystem}
\label{def:direction-frame-selected-channel-system}
Fix a horizon $h$. A selected channel transport system is
\emph{direction-frame isotropic} on a channel fiber $F_h$ if:
\begin{enumerate}[label=(\roman*), leftmargin=2em, itemsep=0.2em]
\item the fiber $F_h$ carries a tight isotropic direction frame in the sense of
      Definition~\ref{def:tight-isotropic-direction-frame};
\item the quadratic transport operator of the selected one-step transport base
      is direction-adapted to that frame in the sense of
      Definition~\ref{def:direction-adapted-quadratic-transport};
\item the selected quadratic transport operator is equivariant for the same
      orthogonal action.
\end{enumerate}
\end{definition}

\begin{corollary}[Realized direction-frame quadratic transport]\lean{selected_channel_scalarization}
\label{cor:realized-direction-frame-quadratic-transport}
Fix a horizon $h$ and let $\mathfrak R_{\mathrm{adm}}(h)$ be a transport-order
disciplined realized class. Suppose a selected channel fiber $F_h$ carries a
direction-frame selected channel system with forced order-one transport base
$B_h$. Then there exists a unique scalar
\[
\kappa_h\ge 0
\]
such that
\[
B_h^*B_h=\kappa_h\,I_{F_h}.
\]
\end{corollary}

\begin{proof}
The hypotheses place $B_h^*B_h$ inside the scope of
Corollary~\ref{cor:direction-frame-quadratic-coefficient}, which yields the
stated scalarization.
\end{proof}

\begin{corollary}[Realized orbit-generated isotropic direction frame]\lean{selected_channel_orbit_direction_frame}
\label{cor:realized-orbit-direction-frame}
Fix a horizon $h$ and let $\mathfrak R_{\mathrm{adm}}(h)$ be a transport-order
disciplined realized class. Suppose a selected channel fiber $F_h$ satisfies:
\begin{enumerate}[label=(\roman*), leftmargin=2em, itemsep=0.2em]
\item there is an orthogonal action
      \[
      U_h:K_h\to O(F_h)
      \]
      making $F_h$ irreducible as a real $K_h$-module;
\item there is a rank-one orthogonal projection
      \[
      \Pi_{0,h}:F_h\to F_h;
      \]
\item the selected direction family is the orbit of $\Pi_{0,h}$ under
      conjugation by $K_h$.
\end{enumerate}
Then the selected direction family is a tight isotropic direction frame on
$F_h$.
\end{corollary}

\begin{proof}
This is exactly Theorem~\ref{thm:orbit-generated-direction-frame} applied to
the selected channel data.
\end{proof}

\begin{definition}[No-preferred-direction selected channel system]\lean{NoPreferredDirectionSelectedChannelSystem}
\label{def:no-preferred-direction-selected-channel-system}
Fix a horizon $h$. A selected channel transport system is \emph{without
preferred direction at quadratic order} on a channel fiber $F_h$ if the
selected one-step transport base carries a quadratic transport rigidity package
in the sense of Definition~\ref{def:quadratic-transport-rigidity-package}.
\end{definition}

\begin{corollary}[Realized no-preferred-direction quadratic transport]\lean{selected_channel_scalarization}
\label{cor:realized-no-preferred-direction-quadratic-transport}
Fix a horizon $h$ and let $\mathfrak R_{\mathrm{adm}}(h)$ be a transport-order
disciplined realized class. Suppose a selected channel fiber $F_h$ carries a
no-preferred-direction selected channel system with forced order-one transport
base $B_h$. Then there exists a unique scalar
\[
\kappa_h\ge 0
\]
such that
\[
B_h^*B_h=\kappa_h\,I_{F_h}.
\]
\end{corollary}

\begin{proof}
The stated data place $B_h$ inside the scope of
Corollary~\ref{cor:no-preferred-direction-quadratic-transport}, which gives
the claimed scalarization.
\end{proof}

\begin{definition}[Quadratically rigid selected channel system]\lean{QuadraticallyRigidSelectedChannelSystem}
\label{def:quadratically-rigid-selected-channel-system}
Fix a horizon $h$. A selected channel transport system is
\emph{quadratically rigid} on a channel fiber $F_h$ if its selected one-step
transport base carries any one of the rigidity packages listed in
Definition~\ref{def:quadratic-transport-rigidity-package}; concretely, this
means one of:
\begin{enumerate}[label=(\roman*), leftmargin=2em, itemsep=0.2em]
\item isotropic selected channel data;
\item sphere-transitive selected channel data;
\item full orthogonal selected channel data;
\item metric-isotropic selected channel data;
\item direction-frame selected channel data.
\end{enumerate}
\end{definition}

\begin{corollary}[Realized quadratic transport rigidity]\lean{selected_channel_scalarization}
\label{cor:realized-quadratic-transport-rigidity}
Fix a horizon $h$ and let $\mathfrak R_{\mathrm{adm}}(h)$ be a transport-order
disciplined realized class. Suppose a selected channel fiber $F_h$ carries a
quadratically rigid selected channel system with forced order-one transport
base $B_h$. Then there exists a unique scalar
\[
\kappa_h\ge 0
\]
such that
\[
B_h^*B_h=\kappa_h\,I_{F_h}.
\]
\end{corollary}

\begin{proof}
The stated data place $B_h$ inside the scope of
Theorem~\ref{thm:quadratic-transport-rigidity}, which gives the claimed
scalarization.
\end{proof}

\begin{corollary}[Realized isotropic recursion channel]\lean{selected_channel_scalar_recursive}
\label{cor:realized-isotropic-recursion-channel}
Fix a horizon $h$ and let $\mathfrak R_{\mathrm{adm}}(h)$ be a transport-order
disciplined realized class. Suppose a selected channel fiber $F_h$ satisfies
the hypotheses of Corollaries~\ref{cor:realized-recursion-forcing} and
\ref{cor:realized-isotropic-quadratic-transport}. Then every bi-infinite
selected channel transport sequence
\[
x=(x_n)_{n\in\bbZ}
\]
satisfies the forced recursion law
\[
x_{n+1}=B_hx_n-x_{n-1},
\]
and the associated second-order quadratic transport data are determined by the
single scalar $\kappa_h$.
\end{corollary}

\begin{proof}
Corollary~\ref{cor:realized-recursion-forcing} gives the recursion law, while
Corollary~\ref{cor:realized-isotropic-quadratic-transport} gives the scalar
description of the quadratic transport operator. Combining them yields the
claimed isotropic recursion channel.
\end{proof}

\begin{corollary}[Realized sphere-transitive recursion channel]\lean{selected_channel_scalar_recursive}
\label{cor:realized-sphere-transitive-recursion-channel}
Fix a horizon $h$ and let $\mathfrak R_{\mathrm{adm}}(h)$ be a transport-order
disciplined realized class. Suppose a selected channel fiber $F_h$ satisfies
the hypotheses of Corollaries~\ref{cor:realized-recursion-forcing} and
\ref{cor:realized-sphere-transitive-quadratic-transport}. Then every
bi-infinite selected channel transport sequence
\[
x=(x_n)_{n\in\bbZ}
\]
satisfies the forced recursion law
\[
x_{n+1}=B_hx_n-x_{n-1},
\]
and the associated second-order quadratic transport data are determined by the
single scalar $\kappa_h$.
\end{corollary}

\begin{proof}
Corollary~\ref{cor:realized-recursion-forcing} gives the recursion law, while
Corollary~\ref{cor:realized-sphere-transitive-quadratic-transport} gives the
scalar description of the quadratic transport operator. Combining them yields
the claimed sphere-transitive recursion channel.
\end{proof}

\begin{corollary}[Realized full orthogonal recursion channel]\lean{selected_channel_scalar_recursive}
\label{cor:realized-full-orthogonal-recursion-channel}
Fix a horizon $h$ and let $\mathfrak R_{\mathrm{adm}}(h)$ be a transport-order
disciplined realized class. Suppose a selected channel fiber $F_h$ satisfies
the hypotheses of Corollaries~\ref{cor:realized-recursion-forcing} and
\ref{cor:realized-full-orthogonal-quadratic-transport}. Then every bi-infinite
selected channel transport sequence
\[
x=(x_n)_{n\in\bbZ}
\]
satisfies the forced recursion law
\[
x_{n+1}=B_hx_n-x_{n-1},
\]
and the associated second-order quadratic transport data are determined by the
single scalar $\kappa_h$.
\end{corollary}

\begin{proof}
Corollary~\ref{cor:realized-recursion-forcing} gives the recursion law, while
Corollary~\ref{cor:realized-full-orthogonal-quadratic-transport} gives the
scalar description of the quadratic transport operator. Combining them yields
the claimed full orthogonal recursion channel.
\end{proof}

\begin{corollary}[Realized metric-isotropic recursion channel]\lean{selected_channel_scalar_recursive}
\label{cor:realized-metric-isotropic-recursion-channel}
Fix a horizon $h$ and let $\mathfrak R_{\mathrm{adm}}(h)$ be a transport-order
disciplined realized class. Suppose a selected channel fiber $F_h$ satisfies
the hypotheses of Corollaries~\ref{cor:realized-recursion-forcing} and
\ref{cor:realized-metric-isotropic-quadratic-transport}. Then every bi-infinite
selected channel transport sequence
\[
x=(x_n)_{n\in\bbZ}
\]
satisfies the forced recursion law
\[
x_{n+1}=B_hx_n-x_{n-1},
\]
and the associated second-order quadratic transport data are determined by the
single scalar $\kappa_h$.
\end{corollary}

\begin{proof}
Corollary~\ref{cor:realized-recursion-forcing} gives the recursion law, while
Corollary~\ref{cor:realized-metric-isotropic-quadratic-transport} gives the
scalar description of the quadratic transport operator. Combining them yields
the claimed metric-isotropic recursion channel.
\end{proof}

\begin{corollary}[Realized direction-frame recursion channel]\lean{selected_channel_scalar_recursive}
\label{cor:realized-direction-frame-recursion-channel}
Fix a horizon $h$ and let $\mathfrak R_{\mathrm{adm}}(h)$ be a transport-order
disciplined realized class. Suppose a selected channel fiber $F_h$ satisfies
the hypotheses of Corollaries~\ref{cor:realized-recursion-forcing} and
\ref{cor:realized-direction-frame-quadratic-transport}. Then every bi-infinite
selected channel transport sequence
\[
x=(x_n)_{n\in\bbZ}
\]
satisfies the forced recursion law
\[
x_{n+1}=B_hx_n-x_{n-1},
\]
and the associated second-order quadratic transport data are determined by the
single scalar $\kappa_h$.
\end{corollary}

\begin{proof}
Corollary~\ref{cor:realized-recursion-forcing} gives the recursion law, while
Corollary~\ref{cor:realized-direction-frame-quadratic-transport} gives the
scalar description of the quadratic transport operator. Combining them yields
the claimed direction-frame recursion channel.
\end{proof}

\begin{corollary}[Realized quadratically rigid recursion channel]\lean{selected_channel_scalar_recursive}
\label{cor:realized-quadratically-rigid-recursion-channel}
Fix a horizon $h$ and let $\mathfrak R_{\mathrm{adm}}(h)$ be a transport-order
disciplined realized class. Suppose a selected channel fiber $F_h$ satisfies
the hypotheses of Corollaries~\ref{cor:realized-recursion-forcing} and
\ref{cor:realized-quadratic-transport-rigidity}. Then every bi-infinite
selected channel transport sequence
\[
x=(x_n)_{n\in\bbZ}
\]
satisfies the forced recursion law
\[
x_{n+1}=B_hx_n-x_{n-1},
\]
and the associated second-order quadratic transport data are determined by the
single scalar $\kappa_h$.
\end{corollary}

\begin{proof}
Corollary~\ref{cor:realized-recursion-forcing} gives the recursion law, while
Corollary~\ref{cor:realized-quadratic-transport-rigidity} gives the scalar
description of the quadratic transport operator. Combining them yields the
claimed quadratically rigid recursion channel.
\end{proof}

\begin{corollary}[Quadratically rigid order-one channels are scalar-recursive]\lean{selected_channel_scalar_recursive}
\label{cor:quadratically-rigid-yields-scalar-recursive}
Fix a horizon $h$ and let $\mathfrak R_{\mathrm{adm}}(h)$ be a transport-order
disciplined realized class. Suppose a selected channel fiber $F_h$ satisfies
the hypotheses of
Corollary~\ref{cor:realized-quadratically-rigid-recursion-channel}. Then that
selected channel system is scalar-recursive in the sense of
Definition~\ref{def:scalar-recursive-selected-channel-system}.
\end{corollary}

\begin{proof}
Corollary~\ref{cor:realized-quadratically-rigid-recursion-channel} gives both
the recursion law and the unique scalar quadratic transport coefficient.
Therefore the selected channel system satisfies
Definition~\ref{def:scalar-recursive-selected-channel-system}.
\end{proof}

\begin{corollary}[Realized no-preferred-direction recursion channel]\lean{selected_channel_scalar_recursive}
\label{cor:realized-no-preferred-direction-recursion-channel}
Fix a horizon $h$ and let $\mathfrak R_{\mathrm{adm}}(h)$ be a transport-order
disciplined realized class. Suppose a selected channel fiber $F_h$ satisfies
the hypotheses of Corollaries~\ref{cor:realized-recursion-forcing} and
\ref{cor:realized-no-preferred-direction-quadratic-transport}. Then every
bi-infinite selected channel transport sequence
\[
x=(x_n)_{n\in\bbZ}
\]
satisfies the forced recursion law
\[
x_{n+1}=B_hx_n-x_{n-1},
\]
and the associated second-order quadratic transport data are determined by the
single scalar $\kappa_h$.
\end{corollary}

\begin{proof}
Corollary~\ref{cor:realized-recursion-forcing} gives the recursion law, while
Corollary~\ref{cor:realized-no-preferred-direction-quadratic-transport} gives
the scalar description of the quadratic transport operator. Combining them
yields the claimed no-preferred-direction recursion channel.
\end{proof}

\begin{definition}[Transport-generated selected channel system]\lean{TransportGeneratedSelectedChannelSystem}
\label{def:transport-generated-selected-channel-system}
Fix a horizon $h$. A selected channel transport system is
\emph{transport-generated} on a channel fiber $F_h$ if:
\begin{enumerate}[label=(\roman*), leftmargin=2em, itemsep=0.2em]
\item the admissible transport bases on $F_h$ satisfy the hypotheses of
      Theorem~\ref{thm:order-one-forcing-criterion};
\item the unique minimal-order admissible equivalence class is represented by a
      bounded transport base
      \[
      B_h:F_h\to F_h;
      \]
\item the selected channel transport system is without preferred direction at
      quadratic order in the sense of
      Definition~\ref{def:no-preferred-direction-selected-channel-system}.
\end{enumerate}
\end{definition}

\begin{corollary}[Realized transport generation forcing]\lean{selected_channel_transport_generation}
\label{cor:realized-transport-generation-forcing}
Fix a horizon $h$ and let $\mathfrak R_{\mathrm{adm}}(h)$ be a transport-order
disciplined realized class. Suppose a selected channel fiber $F_h$ carries a
transport-generated selected channel system with distinguished transport base
$B_h$. Then:
\begin{enumerate}[label=(\roman*), leftmargin=2em, itemsep=0.2em]
\item there exists a unique scalar
      \[
      \kappa_h\ge 0
      \]
      such that
      \[
      B_h^*B_h=\kappa_h\,I_{F_h};
      \]
\item every bi-infinite selected channel transport sequence
      \[
      x=(x_n)_{n\in\bbZ}
      \]
      satisfies the forced recursion law
      \[
      \mathcal R_{B_h}x=0,
      \qquad\text{equivalently}\qquad
      x_{n+1}=B_hx_n-x_{n-1}.
      \]
\end{enumerate}
In particular, the selected channel carries a forced scalar quadratic
recursion datum.
\end{corollary}

\begin{proof}
The stated data place the selected transport channel inside the scope of
Theorem~\ref{thm:transport-generation-forces-scalar-recursion}. Applying that
theorem to the distinguished transport base $B_h$ gives both the unique scalar
coefficient and the forced recursion law.
\end{proof}

\begin{definition}[Scalar-recursive selected channel system]\lean{ScalarRecursiveSelectedChannelSystem}
\label{def:scalar-recursive-selected-channel-system}
Fix a horizon $h$. A selected channel transport system on a channel fiber
$F_h$ with distinguished transport base
\[
B_h:F_h\to F_h
\]
is \emph{scalar-recursive} if:
\begin{enumerate}[label=(\roman*), leftmargin=2em, itemsep=0.2em]
\item there exists a unique scalar
      \[
      \kappa_h\ge 0
      \]
      such that
      \[
      B_h^*B_h=\kappa_h\,I_{F_h};
      \]
\item every bi-infinite selected channel transport sequence
      \[
      x=(x_n)_{n\in\bbZ}
      \]
      satisfies the recursion law
      \[
      x_{n+1}=B_hx_n-x_{n-1}.
      \]
\end{enumerate}
\end{definition}

\begin{definition}[Transport-rigid realized class]\lean{TransportRigidRealizedClass}
\label{def:transport-rigid-realized-class}
Fix a horizon $h$. An admissible realized class
$\mathfrak R_{\mathrm{adm}}(h)$ is \emph{transport-rigid on the selected
channels} if:
\begin{enumerate}[label=(\roman*), leftmargin=2em, itemsep=0.2em]
\item $\mathfrak R_{\mathrm{adm}}(h)$ is transport-order disciplined in the
      sense of
      Definition~\ref{def:transport-order-disciplined-realized-class};
\item on every relevant selected channel fiber $F_h$, the admissible transport
      bases satisfy the hypotheses of
      Theorem~\ref{thm:order-one-forcing-criterion} and the forced one-step
      transport class is represented by a bounded transport base
      \[
      B_h:F_h\to F_h;
      \]
\item every such selected channel system is quadratically rigid in the sense
      of
      Definition~\ref{def:quadratically-rigid-selected-channel-system}.
\end{enumerate}
\end{definition}

\begin{corollary}[Quadratically rigid order-one classes are transport-rigid]\lean{quadratically_rigid_implies_transport_rigid}
\label{cor:quadratically-rigid-implies-transport-rigid}
Fix a horizon $h$ and let $\mathfrak R_{\mathrm{adm}}(h)$ be an admissible
realized class such that:
\begin{enumerate}[label=(\roman*), leftmargin=2em, itemsep=0.2em]
\item $\mathfrak R_{\mathrm{adm}}(h)$ is transport-order disciplined in the
      sense of
      Definition~\ref{def:transport-order-disciplined-realized-class};
\item on every relevant selected channel fiber $F_h$, the admissible transport
      bases satisfy the hypotheses of
      Theorem~\ref{thm:order-one-forcing-criterion} and the forced one-step
      transport class is represented by a bounded transport base
      \[
      B_h:F_h\to F_h;
      \]
\item every relevant selected channel system is quadratically rigid in the
      sense of
      Definition~\ref{def:quadratically-rigid-selected-channel-system}.
\end{enumerate}
Then $\mathfrak R_{\mathrm{adm}}(h)$ is transport-rigid on the selected
channels in the sense of
Definition~\ref{def:transport-rigid-realized-class}.
\end{corollary}

\begin{proof}
Items (i) and (ii) are exactly clauses (i) and (ii) of
Definition~\ref{def:transport-rigid-realized-class}, and item (iii) is clause
(iii). Therefore the class is transport-rigid.
\end{proof}

\begin{corollary}[No-preferred-direction order-one classes are transport-rigid]\lean{no_preferred_direction_implies_transport_rigid}
\label{cor:no-preferred-direction-implies-transport-rigid}
Fix a horizon $h$ and let $\mathfrak R_{\mathrm{adm}}(h)$ be an admissible
realized class such that:
\begin{enumerate}[label=(\roman*), leftmargin=2em, itemsep=0.2em]
\item $\mathfrak R_{\mathrm{adm}}(h)$ is transport-order disciplined in the
      sense of
      Definition~\ref{def:transport-order-disciplined-realized-class};
\item on every relevant selected channel fiber $F_h$, the admissible transport
      bases satisfy the hypotheses of
      Theorem~\ref{thm:order-one-forcing-criterion} and the forced one-step
      transport class is represented by a bounded transport base
      \[
      B_h:F_h\to F_h;
      \]
\item every relevant selected channel system is without preferred direction at
      quadratic order in the sense of
      Definition~\ref{def:no-preferred-direction-selected-channel-system}.
\end{enumerate}
Then $\mathfrak R_{\mathrm{adm}}(h)$ is transport-rigid on the selected
channels in the sense of
Definition~\ref{def:transport-rigid-realized-class}.
\end{corollary}

\begin{proof}
Items (i) and (ii) together with item (iii) place every relevant selected
channel inside the scope of
Definition~\ref{def:quadratically-rigid-selected-channel-system}. Therefore
Corollary~\ref{cor:quadratically-rigid-implies-transport-rigid} applies.
\end{proof}

\begin{corollary}[Intrinsic symmetry/order hypotheses imply governing-datum rigidity]\lean{intrinsic_package_implies_governing_datum_rigid}
\label{cor:intrinsic-package-implies-governing-datum-rigid}
Fix a horizon $h$ and let $\mathfrak R_{\mathrm{adm}}(h)$ be an admissible
realized class such that:
\begin{enumerate}[label=(\roman*), leftmargin=2em, itemsep=0.2em]
\item $\mathfrak R_{\mathrm{adm}}(h)$ is complex-type with scalar multiplicity
      commutant in the sense of
      Definition~\ref{def:complex-type-realized-class};
\item $\mathfrak R_{\mathrm{adm}}(h)$ is observable-irreducible on the
      selected channels in the sense of
      Definition~\ref{def:observable-irreducible-realized-class};
\item the admissible reduced selection data modulo horizon-preserving
      equivalence consist of a single equivalence class;
\item $\mathfrak R_{\mathrm{adm}}(h)$ is transport-order disciplined in the
      sense of
      Definition~\ref{def:transport-order-disciplined-realized-class};
\item on every relevant selected channel fiber $F_h$, the admissible transport
      bases satisfy the hypotheses of
      Theorem~\ref{thm:order-one-forcing-criterion} and the forced one-step
      transport class is represented by a bounded transport base
      \[
      B_h:F_h\to F_h;
      \]
\item every relevant selected channel system is without preferred direction at
      quadratic order in the sense of
      Definition~\ref{def:no-preferred-direction-selected-channel-system}.
\end{enumerate}
Then $\mathfrak R_{\mathrm{adm}}(h)$ is governing-datum rigid on the selected
channels in the sense of
Definition~\ref{def:governing-datum-rigid-realized-class}.
\end{corollary}

\begin{proof}
By items (i)--(iii) and
Corollary~\ref{cor:scalar-commutant-complex-type-implies-observable-rigid},
$\mathfrak R_{\mathrm{adm}}(h)$ is observable-rigid on the selected channels.
By items (iv)--(vi) and
Corollary~\ref{cor:no-preferred-direction-implies-transport-rigid},
$\mathfrak R_{\mathrm{adm}}(h)$ is transport-rigid on the selected channels.
Therefore both clauses of
Definition~\ref{def:governing-datum-rigid-realized-class} hold.
\end{proof}

\begin{corollary}[Transport-rigid classes are scalar-recursive]\lean{transport_rigid_implies_scalar_recursive}
\label{cor:transport-rigid-implies-scalar-recursive}
Fix a horizon $h$ and let $\mathfrak R_{\mathrm{adm}}(h)$ be a transport-rigid
realized class on the selected channels. Then every relevant selected channel
fiber is scalar-recursive in the sense of
Definition~\ref{def:scalar-recursive-selected-channel-system}.
\end{corollary}

(Proof in Appendix~\ref{sec:appendix-dynamics}.)

\begin{corollary}[Transport generation yields a scalar-recursive channel]\lean{transport_generation_yields_scalar_recursive}
\label{cor:transport-generation-yields-scalar-recursive}
Fix a horizon $h$ and let $\mathfrak R_{\mathrm{adm}}(h)$ be a transport-order
disciplined realized class. Suppose a selected channel fiber $F_h$ carries a
transport-generated selected channel system with distinguished transport base
$B_h$. Then that selected channel system is scalar-recursive in the sense of
Definition~\ref{def:scalar-recursive-selected-channel-system}.
\end{corollary}

(Proof in Appendix~\ref{sec:appendix-dynamics}.)

\begin{corollary}[Transport-generated classes are transport-rigid]\lean{transport_generated_implies_transport_rigid}
\label{cor:transport-generated-implies-transport-rigid}
Fix a horizon $h$ and let $\mathfrak R_{\mathrm{adm}}(h)$ be a transport-order
disciplined realized class. Suppose every relevant selected channel fiber
carries a transport-generated selected channel system in the sense of
Definition~\ref{def:transport-generated-selected-channel-system}. Then
$\mathfrak R_{\mathrm{adm}}(h)$ is transport-rigid on the selected channels in
the sense of Definition~\ref{def:transport-rigid-realized-class}.
\end{corollary}

(Proof in Appendix~\ref{sec:appendix-dynamics}.)

\begin{definition}[Observable-rigidity generated realized class]\lean{ObservableRigidityGeneratedRealizedClass}
\label{def:observable-rigidity-generated-realized-class}
Fix a horizon $h$. An admissible realized class
$\mathfrak R_{\mathrm{adm}}(h)$ is \emph{observable-rigidity generated on the
selected channels} if, on every relevant selected-law channel, at least one of
the following realized witness packages holds:
\begin{enumerate}[label=(\roman*), leftmargin=2em, itemsep=0.2em]
\item a phase-generated realized-class package in the sense of
      Definition~\ref{def:phase-generated-realized-class}, together with
      observable irreducibility on the selected channels;
\item an observable-complete realized-class package in the sense of
      Corollary~\ref{cor:observable-complete-implies-observable-rigid};
\item a finite-observable-rigidity realized-class package in the sense of
      Corollary~\ref{cor:finite-observable-rigidity-implies-observable-rigid};
\item an adjoint-closed real-type realized-class package in the sense of
      Corollary~\ref{cor:adjoint-closed-real-type-implies-observable-rigid}.
\item an observable-orbit generated realized-class package in the sense of
      Definition~\ref{def:observable-orbit-generated-realized-class}.
\end{enumerate}
\end{definition}

\begin{corollary}[Observable-rigidity generated classes are observable-rigid]\lean{observable_rigidity_generated_implies_rigid}
\label{cor:observable-rigidity-generated-implies-rigid}
Fix a horizon $h$ and let $\mathfrak R_{\mathrm{adm}}(h)$ be an
observable-rigidity generated realized class on the selected channels. Then
$\mathfrak R_{\mathrm{adm}}(h)$ is observable-rigid on the selected channels in
the sense of Definition~\ref{def:observable-rigid-realized-class}.
\end{corollary}

\begin{proof}
Apply, channel by channel, the corresponding route from
Corollary~\ref{cor:generated-observable-implies-rigid},
Corollary~\ref{cor:observable-complete-implies-observable-rigid},
Corollary~\ref{cor:finite-observable-rigidity-implies-observable-rigid}, or
Corollary~\ref{cor:adjoint-closed-real-type-implies-observable-rigid}, or
Corollary~\ref{cor:observable-orbit-generated-implies-observable-rigid},
according to which witness package is available.
\end{proof}

\begin{definition}[Transport-rigidity generated realized class]\lean{TransportRigidityGeneratedRealizedClass}
\label{def:transport-rigidity-generated-realized-class}
Fix a horizon $h$. An admissible realized class
$\mathfrak R_{\mathrm{adm}}(h)$ is \emph{transport-rigidity generated on the
selected channels} if, on every relevant selected channel fiber, at least one
of the following realized witness packages holds:
\begin{enumerate}[label=(\roman*), leftmargin=2em, itemsep=0.2em]
\item a transport-generated selected channel system in the sense of
      Definition~\ref{def:transport-generated-selected-channel-system};
\item a no-preferred-direction order-one package in the sense of
      Corollary~\ref{cor:no-preferred-direction-implies-transport-rigid};
\item a quadratically rigid order-one package in the sense of
      Corollary~\ref{cor:quadratically-rigid-implies-transport-rigid}.
\end{enumerate}
\end{definition}

\begin{corollary}[Transport-rigidity generated classes are transport-rigid]\lean{transport_rigidity_generated_implies_rigid}
\label{cor:transport-rigidity-generated-implies-rigid}
Fix a horizon $h$ and let $\mathfrak R_{\mathrm{adm}}(h)$ be a
transport-rigidity generated realized class on the selected channels. Then
$\mathfrak R_{\mathrm{adm}}(h)$ is transport-rigid on the selected channels in
the sense of Definition~\ref{def:transport-rigid-realized-class}.
\end{corollary}

\begin{proof}
Apply, channel by channel, the corresponding route from
Corollary~\ref{cor:transport-generated-implies-transport-rigid},
Corollary~\ref{cor:no-preferred-direction-implies-transport-rigid}, or
Corollary~\ref{cor:quadratically-rigid-implies-transport-rigid}, according to
which witness package is available.
\end{proof}

\begin{definition}[Rigidity-generated realized class]\lean{RigidityGeneratedRealizedClass}
\label{def:rigidity-generated-realized-class}
Fix a horizon $h$. An admissible realized class
$\mathfrak R_{\mathrm{adm}}(h)$ is \emph{rigidity-generated on the selected
channels} if:
\begin{enumerate}[label=(\roman*), leftmargin=2em, itemsep=0.2em]
\item it is observable-rigidity generated on the selected channels in the
      sense of
      Definition~\ref{def:observable-rigidity-generated-realized-class};
\item it is transport-rigidity generated on the selected channels in the sense
      of
      Definition~\ref{def:transport-rigidity-generated-realized-class}.
\end{enumerate}
\end{definition}

\begin{corollary}[Rigidity-generated classes are governing-datum rigid]\lean{rigidity_generated_implies_governing_datum_rigid}
\label{cor:rigidity-generated-implies-governing-datum-rigid}
Fix a horizon $h$ and let $\mathfrak R_{\mathrm{adm}}(h)$ be a
rigidity-generated realized class on the selected channels. Then
$\mathfrak R_{\mathrm{adm}}(h)$ is governing-datum rigid on the selected
channels in the sense of
Definition~\ref{def:governing-datum-rigid-realized-class}.
\end{corollary}

\begin{proof}
By Corollary~\ref{cor:observable-rigidity-generated-implies-rigid},
$\mathfrak R_{\mathrm{adm}}(h)$ is observable-rigid on the selected channels.
By Corollary~\ref{cor:transport-rigidity-generated-implies-rigid},
$\mathfrak R_{\mathrm{adm}}(h)$ is transport-rigid on the selected channels.
Therefore both clauses of
Definition~\ref{def:governing-datum-rigid-realized-class} hold.
\end{proof}

\begin{definition}[Governing-datum rigid realized class]\lean{GoverningDatumRigidRealizedClass}
\label{def:governing-datum-rigid-realized-class}
Fix a horizon $h$. An admissible realized class
$\mathfrak R_{\mathrm{adm}}(h)$ is \emph{governing-datum rigid on the selected
channels} if:
\begin{enumerate}[label=(\roman*), leftmargin=2em, itemsep=0.2em]
\item $\mathfrak R_{\mathrm{adm}}(h)$ is observable-rigid on the selected
      channels in the sense of
      Definition~\ref{def:observable-rigid-realized-class};
\item $\mathfrak R_{\mathrm{adm}}(h)$ is transport-rigid on the selected
      channels in the sense of
      Definition~\ref{def:transport-rigid-realized-class}.
\end{enumerate}
\end{definition}

\begin{definition}[Closure-rigid realized class]\lean{ClosureRigidRealizedClass}
\label{def:closure-rigid-realized-class}
Fix a horizon $h$. An admissible realized class
$\mathfrak R_{\mathrm{adm}}(h)$ is \emph{closure-rigid on the selected
symmetry channels} if, for every realization $R\in\mathfrak R_{\mathrm{adm}}(h)$:
\begin{enumerate}[label=(\roman*), leftmargin=2em, itemsep=0.2em]
\item the observed sector $\CCHobs{h}(R)$ carries a finite unitary symmetry
      action and an isotypic decomposition on the relevant selected-law sector;
\item every selected self-adjoint closure operator on that sector is
      equivariant for the chosen symmetry action;
\item every such selected self-adjoint closure operator is symmetry-scalar on
      every appearing isotypic channel in the sense of
      Definition~\ref{def:symmetry-scalar-closure}.
\end{enumerate}
\end{definition}

\begin{corollary}[Multiplicity-free equivariant closure classes are closure-rigid]\lean{multiplicity_free_equivariant_closure_implies_rigid}
\label{cor:multiplicity-free-equivariant-closure-implies-rigid}
Fix a horizon $h$ and let $\mathfrak R_{\mathrm{adm}}(h)$ be an admissible
realized class such that, for every realization $R\in\mathfrak R_{\mathrm{adm}}(h)$:
\begin{enumerate}[label=(\roman*), leftmargin=2em, itemsep=0.2em]
\item the observed sector $\CCHobs{h}(R)$ carries a finite unitary symmetry
      action and an isotypic decomposition on the relevant selected-law sector;
\item every selected self-adjoint closure operator on that sector is
      equivariant for the chosen symmetry action;
\item the symmetry representation on the relevant selected-law sector is
      multiplicity-free.
\end{enumerate}
Then $\mathfrak R_{\mathrm{adm}}(h)$ is closure-rigid on the selected symmetry
channels in the sense of
Definition~\ref{def:closure-rigid-realized-class}.
\end{corollary}

\begin{proof}
Fix a realization $R\in\mathfrak R_{\mathrm{adm}}(h)$ and a selected
self-adjoint closure operator $A$ on the relevant selected-law sector. By
items (i)--(iii), the operator $A$ lies in the scope of
Theorem~\ref{thm:multiplicity-free-closures-scalar}. Hence $A$ is
symmetry-scalar on every appearing isotypic channel, which is exactly clause
(iii) of Definition~\ref{def:closure-rigid-realized-class}. Clauses (i) and
(ii) are part of the assumptions, so the definition holds for every
realization.
\end{proof}

\begin{definition}[Symmetry-generated realized class]\lean{SymmetryGeneratedRealizedClass}
\label{def:symmetry-generated-realized-class}
Fix a horizon $h$. An admissible realized class
$\mathfrak R_{\mathrm{adm}}(h)$ is \emph{symmetry-generated on the selected
channels} if:
\begin{enumerate}[label=(\roman*), leftmargin=2em, itemsep=0.2em]
\item it is observable-rigidity generated on the selected channels in the
      sense of
      Definition~\ref{def:observable-rigidity-generated-realized-class};
\item it is transport-order disciplined in the sense of
      Definition~\ref{def:transport-order-disciplined-realized-class};
\item on every relevant selected channel fiber, it carries a
      transport-generated selected channel system in the sense of
      Definition~\ref{def:transport-generated-selected-channel-system};
\item for every realization $R\in\mathfrak R_{\mathrm{adm}}(h)$, the observed
      sector $\CCHobs{h}(R)$ carries a finite unitary symmetry action and an
      isotypic decomposition on the relevant selected-law sector;
\item every selected self-adjoint closure operator on that sector is
      equivariant for the chosen symmetry action;
\item the symmetry representation on the relevant selected-law sector is
      multiplicity-free.
\end{enumerate}
\end{definition}

\begin{corollary}[Symmetry-generated classes are governing-datum rigid]\lean{symmetry_generated_implies_governing_datum_rigid}
\label{cor:symmetry-generated-implies-governing-datum-rigid}
Fix a horizon $h$ and let $\mathfrak R_{\mathrm{adm}}(h)$ be a
symmetry-generated realized class on the selected channels. Then
$\mathfrak R_{\mathrm{adm}}(h)$ is governing-datum rigid on the selected
channels in the sense of
Definition~\ref{def:governing-datum-rigid-realized-class}.
\end{corollary}

\begin{proof}
By Corollary~\ref{cor:observable-rigidity-generated-implies-rigid},
$\mathfrak R_{\mathrm{adm}}(h)$ is observable-rigid on the selected channels.
By items (ii)--(iii) of
Definition~\ref{def:symmetry-generated-realized-class} and
Corollary~\ref{cor:transport-generated-implies-transport-rigid},
$\mathfrak R_{\mathrm{adm}}(h)$ is transport-rigid on the selected channels.
Therefore both clauses of
Definition~\ref{def:governing-datum-rigid-realized-class} hold.
\end{proof}

\begin{corollary}[Symmetry-generated classes are closure-rigid]\lean{symmetry_generated_implies_closure_rigid}
\label{cor:symmetry-generated-implies-closure-rigid}
Fix a horizon $h$ and let $\mathfrak R_{\mathrm{adm}}(h)$ be a
symmetry-generated realized class on the selected channels. Then
$\mathfrak R_{\mathrm{adm}}(h)$ is closure-rigid on the selected symmetry
channels in the sense of
Definition~\ref{def:closure-rigid-realized-class}.
\end{corollary}

\begin{proof}
Items (iv)--(vi) of Definition~\ref{def:symmetry-generated-realized-class}
are exactly the hypotheses of
Corollary~\ref{cor:multiplicity-free-equivariant-closure-implies-rigid}.
The conclusion follows directly.
\end{proof}

\begin{definition}[Locally homogeneous realized class]\lean{LocallyHomogeneousRealizedClass}
\label{def:locally-homogeneous-realized-class}
Fix a horizon $h$. An admissible realized class
$\mathfrak R_{\mathrm{adm}}(h)$ is \emph{locally homogeneous on the selected
channels} if:
\begin{enumerate}[label=(\roman*), leftmargin=2em, itemsep=0.2em]
\item it is observable-orbit generated on the selected channels in the sense
      of Definition~\ref{def:observable-orbit-generated-realized-class};
\item it is transport-order disciplined in the sense of
      Definition~\ref{def:transport-order-disciplined-realized-class};
\item on every relevant selected channel fiber $F_h$, there is an orthogonal
      action
      \[
      U_h:K_h\to O(F_h)
      \]
      making $F_h$ irreducible as a real $K_h$-module, together with a
      rank-one orthogonal projection
      \[
      \Pi_{0,h}:F_h\to F_h
      \]
      whose orbit under conjugation by $K_h$ is the selected direction family;
\item on every relevant selected channel fiber, the selected quadratic
      transport operator is direction-adapted to that orbit frame and
      equivariant for the same orthogonal action;
\item on every relevant selected channel fiber, the admissible transport bases
      satisfy the hypotheses of
      Theorem~\ref{thm:order-one-forcing-criterion} and the forced order-one
      class is represented by a bounded transport base;
\item for every realization $R\in\mathfrak R_{\mathrm{adm}}(h)$, the observed
      sector $\CCHobs{h}(R)$ carries a finite unitary symmetry action and an
      isotypic decomposition on the relevant selected-law sector;
\item every selected self-adjoint closure operator on that sector is
      equivariant for the chosen symmetry action;
\item the symmetry representation on the relevant selected-law sector is
      multiplicity-free.
\end{enumerate}
\end{definition}

\begin{corollary}[Observable-orbit orbit-frame multiplicity-free classes are symmetry-generated]\lean{observable_orbit_frame_implies_symmetry_generated}
\label{cor:observable-orbit-frame-implies-symmetry-generated}
Fix a horizon $h$ and let $\mathfrak R_{\mathrm{adm}}(h)$ be an admissible
realized class such that:
\begin{enumerate}[label=(\roman*), leftmargin=2em, itemsep=0.2em]
\item $\mathfrak R_{\mathrm{adm}}(h)$ is observable-orbit generated on the
      selected channels in the sense of
      Definition~\ref{def:observable-orbit-generated-realized-class};
\item $\mathfrak R_{\mathrm{adm}}(h)$ is transport-order disciplined in the
      sense of
      Definition~\ref{def:transport-order-disciplined-realized-class};
\item on every relevant selected channel fiber $F_h$, there is an orthogonal
      action
      \[
      U_h:K_h\to O(F_h)
      \]
      making $F_h$ irreducible as a real $K_h$-module, together with a
      rank-one orthogonal projection
      \[
      \Pi_{0,h}:F_h\to F_h
      \]
      whose orbit under conjugation by $K_h$ is the selected direction family;
\item on every relevant selected channel fiber, the selected quadratic
      transport operator is direction-adapted to that orbit frame and
      equivariant for the same orthogonal action;
\item on every relevant selected channel fiber, the admissible transport bases
      satisfy the hypotheses of
      Theorem~\ref{thm:order-one-forcing-criterion} and the forced order-one
      class is represented by a bounded transport base;
\item for every realization $R\in\mathfrak R_{\mathrm{adm}}(h)$, the observed
      sector $\CCHobs{h}(R)$ carries a finite unitary symmetry action and an
      isotypic decomposition on the relevant selected-law sector;
\item every selected self-adjoint closure operator on that sector is
      equivariant for the chosen symmetry action;
\item the symmetry representation on the relevant selected-law sector is
      multiplicity-free.
\end{enumerate}
Then $\mathfrak R_{\mathrm{adm}}(h)$ is symmetry-generated on the selected
channels in the sense of
Definition~\ref{def:symmetry-generated-realized-class}.
\end{corollary}

(Proof in Appendix~\ref{sec:appendix-dynamics}.)

\begin{corollary}[Locally homogeneous classes are symmetry-generated]\lean{locally_homogeneous_implies_symmetry_generated}
\label{cor:locally-homogeneous-implies-symmetry-generated}
Fix a horizon $h$ and let $\mathfrak R_{\mathrm{adm}}(h)$ be a locally
homogeneous realized class on the selected channels. Then
$\mathfrak R_{\mathrm{adm}}(h)$ is symmetry-generated on the selected channels
in the sense of Definition~\ref{def:symmetry-generated-realized-class}.
\end{corollary}

\begin{proof}
The hypotheses of
Definition~\ref{def:locally-homogeneous-realized-class} are exactly the
hypotheses of
Corollary~\ref{cor:observable-orbit-frame-implies-symmetry-generated}. The
conclusion follows directly.
\end{proof}

\begin{definition}[Local internal datum on selected channels]\lean{LocalInternalDatumSelectedChannels}
\label{def:local-internal-datum-selected-channels}
Fix a horizon $h$ and let $\mathfrak R_{\mathrm{adm}}(h)$ be a governing-datum
rigid realized class. For each relevant selected channel
$\rho$ with fiber $F_{h,\rho}$ and distinguished transport base $B_{h,\rho}$,
the \emph{local internal datum} at horizon $h$ is the channelwise package
\[
\mathfrak g_{h,\rho}:=\bigl([J_{h,\rho}],\,\kappa_{h,\rho},\,[C_{h,\rho}]\bigr),
\]
where:
\begin{enumerate}[label=(\roman*), leftmargin=2em, itemsep=0.2em]
\item $[J_{h,\rho}]$ is the forced phase-carrier class modulo
      phase-equivalence;
\item $\kappa_{h,\rho}\ge 0$ is the unique scalar for which
      \[
      B_{h,\rho}^*B_{h,\rho}=\kappa_{h,\rho}\,I_{F_{h,\rho}};
      \]
\item $[C_{h,\rho}]$ denotes the scalar closure class on the corresponding
      multiplicity space for every selected self-adjoint closure operator on
      that channel.
\end{enumerate}
\end{definition}

\begin{theorem}[Intrinsic fixed-horizon governing datum forcing]\lean{realized_local_internal_datum_forcing}
\label{cor:realized-local-internal-datum-forcing}
Fix a horizon $h$ and let $\mathfrak R_{\mathrm{adm}}(h)$ be a governing-datum
rigid realized class. Then on every relevant selected
channel:
\begin{enumerate}[label=(\roman*), leftmargin=2em, itemsep=0.2em]
\item the selected law carries a forced phase-carrier class modulo
      phase-equivalence;
\item there exists a unique scalar
      \[
      \kappa_{h,\rho}\ge 0
      \]
      such that the distinguished transport base satisfies
      \[
      B_{h,\rho}^*B_{h,\rho}=\kappa_{h,\rho}\,I_{F_{h,\rho}},
      \]
      and every bi-infinite selected channel transport sequence satisfies the
      forced recursion law
      \[
      x_{n+1}=B_{h,\rho}x_n-x_{n-1};
      \]
\item every selected self-adjoint closure operator on that channel is scalar on
      the corresponding multiplicity space.
\end{enumerate}
In particular, the selected law carries a unique local internal datum in the
sense of Definition~\ref{def:local-internal-datum-selected-channels}.
\end{theorem}

(Proof in Appendix~\ref{sec:appendix-dynamics}.)

\begin{definition}[Gauge-compatible local internal datum]\lean{GaugeCompatibleLocalInternalDatum}
\label{def:gauge-compatible-local-internal-datum}
Fix a horizon $h$ and let $\mathfrak R_{\mathrm{adm}}(h)$ be a governing-datum
rigid realized class. Assume the relevant selected channels carry selected
channel transport systems in the sense of
Definition~\ref{def:selected-channel-transport-system}. For each such selected
channel $\rho$, the \emph{gauge-compatible local internal datum} is the
quadruple
\[
\widehat{\mathfrak g}_{h,\rho}
:=
\bigl([J_{h,\rho}],\,\kappa_{h,\rho},\,[C_{h,\rho}],\,[\mathcal U_{h,\rho}]\bigr),
\]
where:
\begin{enumerate}[label=(\roman*), leftmargin=2em, itemsep=0.2em]
\item $\bigl([J_{h,\rho}],\,\kappa_{h,\rho},\,[C_{h,\rho}]\bigr)$ is the local
      internal datum of Definition~\ref{def:local-internal-datum-selected-channels};
\item $[\mathcal U_{h,\rho}]$ is the local unitary gauge-covariance class of
      the selected channel transport system on $\rho$, modulo independent local
      unitary frame changes on the selected channel fibers;
\item if preserved unit volume forms are part of the realized data, then
      $[\mathcal U_{h,\rho}]$ is understood in the reduced unimodular gauge
      class.
\end{enumerate}
\end{definition}

\begin{theorem}[Gauge-compatible fixed-horizon governing datum forcing]\lean{gauge_compatible_local_datum_forcing}
\label{thm:gauge-compatible-local-datum-forcing}
Fix a horizon $h$ and let $\mathfrak R_{\mathrm{adm}}(h)$ be a governing-datum
rigid realized class. Assume the relevant selected channels carry selected
channel transport systems in the sense of
Definition~\ref{def:selected-channel-transport-system}. Then on every relevant
selected channel:
\begin{enumerate}[label=(\roman*), leftmargin=2em, itemsep=0.2em]
\item the selected law carries a unique local internal datum in the sense of
      Definition~\ref{def:local-internal-datum-selected-channels};
\item the selected channel transport system is locally unitary gauge-covariant;
\item if preserved unit volume forms are part of the realized data, the local
      gauge group reduces to the unimodular subgroup.
\end{enumerate}
In particular, the selected law carries a unique gauge-compatible local
internal datum in the sense of
Definition~\ref{def:gauge-compatible-local-internal-datum}.
\end{theorem}

\begin{proof}
Item (i) is exactly
Theorem~\ref{cor:realized-local-internal-datum-forcing}. Item (ii) is
Corollary~\ref{cor:realized-local-gauge-covariance}(i), and item (iii) is
Corollary~\ref{cor:realized-local-gauge-covariance}(ii). These are precisely
the data listed in
Definition~\ref{def:gauge-compatible-local-internal-datum}.
\end{proof}

\begin{corollary}[Intrinsic symmetry/order hypotheses yield gauge-compatible local datum]\lean{intrinsic_package_yields_gauge_compatible_local_datum}
\label{cor:intrinsic-package-yields-gauge-compatible-local-datum}
Fix a horizon $h$ and let $\mathfrak R_{\mathrm{adm}}(h)$ be an admissible
realized class satisfying the hypotheses of
Corollary~\ref{cor:intrinsic-package-implies-governing-datum-rigid}. Assume
further that the relevant selected channels carry selected channel transport
systems in the sense of
Definition~\ref{def:selected-channel-transport-system}. Then the selected law
carries a unique gauge-compatible local internal datum in the sense of
Definition~\ref{def:gauge-compatible-local-internal-datum}.
\end{corollary}

\begin{proof}
By Corollary~\ref{cor:intrinsic-package-implies-governing-datum-rigid},
$\mathfrak R_{\mathrm{adm}}(h)$ is governing-datum rigid on the selected
channels. The conclusion is therefore exactly
Theorem~\ref{thm:gauge-compatible-local-datum-forcing}.
\end{proof}

\paragraph{Boundary of Chapter~13.}
\label{rem:part-two-boundary}
This chapter stops at forcing on realized classes at fixed horizons. The next
question concerns refinement: when selected laws are organized across
successive horizons, how do they recover classical continuum operators, and
under what additional hypotheses is the continuum law itself forced? That
classical-limit bridge is taken up separately in Chapter~9.
